\documentclass[12pt]{amsart} 

\usepackage{amsfonts,amssymb,stmaryrd,amscd,amsmath,latexsym,amsbsy,xcolor,tikz-cd,hyperref}
\numberwithin{equation}{section}
\definecolor{commentred}{RGB}{190,0,0}
\definecolor{additionpurple}{RGB}{120,0,160}
\definecolor{suggestionblue}{RGB}{0,70,180}
\newcommand{\redcomment}[1]{\par\smallskip\noindent{\color{commentred}\textbf{Comment.} #1}\par\smallskip}
\newenvironment{redadd}{\par\smallskip\begingroup\color{commentred}}{\par\smallskip\endgroup}
\newenvironment{purpleadd}{\par\smallskip\begingroup\color{additionpurple}}{\par\smallskip\endgroup}
\newenvironment{blueadd}{\par\smallskip\begingroup\color{suggestionblue}}{\par\smallskip\endgroup}
\hypersetup{hidelinks}

\newtheorem{theorem}{Theorem}[section]
\newtheorem{lemma}[theorem]{Lemma}
\newtheorem{proposition}[theorem]{Proposition}
\newtheorem{corollary}[theorem]{Corollary}
\theoremstyle{definition}
\newtheorem{definition}[theorem]{Definition}

\newtheorem{example}[theorem]{Example}
\newtheorem{question}[theorem]{Question}
\newtheorem{conjecture}[theorem]{Conjecture}
\newtheorem{remark}[theorem]{Remark}
\newtheorem{note}[theorem]{Note}

%%%%%%%%%%%%%%% symbols shorthand %%%%%%%%%%%
%\renewcommand{\text}[1]{\mbox{{\rm #1}}}
\newcommand{\Ext}{\text{Ext}}
\newcommand{\Hh}{\mathcal{H}}
\newcommand{\Tr}{\text{Tr}}%{\text{Tr}\,}
\newcommand{\id}{\text{id}}
\newcommand{\Image}{\text{Im}}
\newcommand{\Ker}{\text{Ker\,}}
\newcommand{\rank}{\text{rank\,}}
\newcommand{\Index}{\text{Index}\,}
\newcommand{\Map}{\text{Map}}
\newcommand{\End}{\text{End}}
\newcommand{\Hom}{\text{Hom}}
\newcommand{\Ad}{\text{Ad}}
\newcommand{\Id}{\text{Id}}
\newcommand{\Aut}{\text{Aut}}
\newcommand{\Sh}{\text{Sh}}
\newcommand{\Ind}{\text{Ind}}
\newcommand{\Rep}{\text{Rep}}
\newcommand{\supp}{\text{supp}\,}
\renewcommand{\span}{\text{span}}
\newcommand{\depth}{\text{depth}}
\newcommand{\St}{\text{St}}
\newcommand{\diag}{\text{diag}}
\newcommand{\comp}{\text{comp}}
\newcommand{\op}{{\text{op}}}
\newcommand{\tr}{\text{tr}\,}


\newcommand{\eps}{\varepsilon}
\newcommand{\mR}{{\mathcal R}}
\newcommand{\bmR}{{\bar \mR}}
\newcommand{\bR}{{\bar R}}
\newcommand{\bV}{{\bar V}}
\newcommand{\bF}{{\bar F}}
\newcommand{\I}{\mathcal{I}}
\renewcommand{\O}{\mathcal{O}}

\newcommand{\G}{\widetilde{G}}

\newcommand{\g}{\mathfrak{g}}
%\newcommand{\kk}{\mathfrak{k}}
\newcommand{\p}{\mathfrak{p}}
\newcommand{\h}{\mathfrak{h}}
\newcommand{\n}{\mathfrak{n}}
\renewcommand{\b}{\mathfrak{b}}
\newcommand{\wg}{\widehat{\mathfrak{g}}}
\newcommand{\wh}{\widehat{\mathfrak{h}}}
\newcommand{\wn}{\widehat{\mathfrak{n}}}
\newcommand{\wb}{\widehat{\mathfrak{b}}}
\newcommand{\ov}{\overline}
\newcommand{\Dim}{{\rm Dim}}
\newcommand{\kk}{{\bold k}}
\newcommand{\actl}{\rightharpoonup}
\newcommand{\actr}{\leftharpoonup}
\newcommand{\lact}{\triangleright}
\newcommand{\ract}{\triangleleft}
\newcommand{\la}{\lambda}
\newcommand{\ra}{\,\rangle}
\newcommand{\flip}{\varsigma}
\newcommand{\nv}{\vec{n}}
\renewcommand{\star}{{\dag}}
%\newcommand{\rtimes}{{>\!\!\!\triangleleft}}
%\newcommand{\ltimes}{{\triangleright\!\!\!<}}
\newcommand{\smrtimes}{>\!\!\triangleleft}
\newcommand{\equiotient}{\Leftrightarrow}
\newcommand{\half}{1/2}%%%%%{\frac{1}{2}}
\newcommand{\own}{\ni}
\newcommand{\1}{_{(1)}}
\newcommand{\2}{_{(2)}}
\newcommand{\3}{_{(3)}}
\newcommand{\4}{_{(4)}}
\newcommand{\5}{_{(5)}}
\newcommand{\6}{_{(6)}}
\renewcommand{\I}{^{(1)}}
\newcommand{\II}{^{(2)}}
\newcommand{\tDelta}{\Delta^{op}}
\newcommand{\tS}{\tilde{S}}
\newcommand{\teps}{\tilde\eps}
\newcommand{\tI}{_{(\tilde{1})}}
\newcommand{\tII}{_{(\tilde{2})}}
\newcommand{\tIII}{_{(\tilde{3})}}
\newcommand{\bTheta}{{\bar \Theta}}
\newcommand{\bs}{\backslash}

\newcommand{\C}{\mathcal{C}}
\newcommand{\ot}{\otimes}
\newcommand{\qexp}{{\text{qexp}}}
\newcommand{\ben}{\begin{enumerate}}
\newcommand{\een}{\end{enumerate}}
\newcommand{\ad}{{\text{Ad}}}
\newcommand{\Rad}{{\text{Rad}}}
\newcommand{\Vect}{{\text{Vect}}}
\newcommand{\Lie}{{\text{Lie}}}
\newcommand{\Cl}{{\text{Cl}}}
\newcommand{\mC}{{\mathcal C}}
\newcommand{\mB}{{\mathcal B}}
\newcommand{\mD}{{\mathcal D}}
\newcommand{\trace}{{\rm tr}}
\newcommand{\CC}{{\mathbb{C}}}
\newcommand{\ZZ}{{\mathbb{Z}}}
\newcommand{\mM}{{\mathcal M}}
\newcommand{\be}{{\bf 1}}

\theoremstyle{plain}

\newtheorem{thm}{Theorem}[section]

\newtheorem*{thrm}{Theorem}

\newtheorem{prop}{Proposition}[section]

\newtheorem{pro}{Proposition}[section]

%\newtheorem{defi}{Definition}

\newtheorem{lem}{Lemma}[section]

\newtheorem{cor}{Corollary}[section]

\newtheorem{eg}{{Example}}

\newtheorem{ex}{Exercise}

\newtheorem*{sol}{Solution}

\newtheorem*{LRrule}{Littlewood-Richardson rule}

\newtheorem*{conj}{Conjecture}



\theoremstyle{definition}

\newtheorem*{defi}{Definition}



\theoremstyle{remark}

\newtheorem{rema}{Remark}[section]





\newcommand{\solu}[1]{\begin{sol}{\bf (\ref{#1})}}

\newcommand{\dd}{\partial}



\newcommand{\gotk}{\mathfrak k}

\def\intersect{\mathop{\bigcap}}
\def\union{\mathop{\bigcup}}
\def\g{\mathfrak{g}}
\def\R{\mathbb{R}}
\def\u{\mathfrak{u}}
\def\su{\mathfrak{su}}
\def\gl{\mathfrak{gl}}
\def\sl{\mathfrak{sl}}
\def\tr{\mathrm{Tr}}
\def\C{\mathcal{C}}
\def\D{\mathcal{D}}
\def\E{\mathcal{E}}
\def\Der{\mathrm{Der}}
\def\Aut{\mathop{\mathrm{Aut}}\nolimits}
\def\h{\mathfrak{h}}
\def\so{\mathfrak{so}}
\def\ad{\mathop{\mathrm{ad}}\nolimits}
\def\Ad{\mathop{\mathrm{Ad}}\nolimits}
\def\Vect{\mathrm{Vect}}
\def\P{\mathcal{P}}
\def\Z{\mathcal{Z}}
\def\Q{\mathbb{Q}}
\def\F{\mathbb{F}}
\def\spn{\mathrm{span}}
\def\H{\mathcal{H}}
\def\rad{\mathop{rad}}
\def\N{\mathcal{N}}
\def\M{\mathcal{M}}
\def\coad{\mathop{coad}}
\def\z{\mathfrak{z}}
\def\a{\mathfrak{a}}
\def\End{\mathrm{End}}
\def\Hom{\mathrm{Hom}}
\def\Ker{\mathrm{Ker}}
\def\Im{\mathop{Im}}
\def\Gr{\mathop{Gr}}
\def\n{\mathfrak{n}}
\def\r{\mathfrak{r}}
\def\soc{\mathop{\mathrm{soc}}\nolimits}
\def\nil{\mathop{\mathrm{nil}}\nolimits}
\def\sp{\mathfrak{sp}}
\def\Dim{{\rm Dim}}
\def\Lie{\mathop{\mathrm{Lie}}\nolimits}
\def\B{\mathcal{B}}
\def\GL{\mathop{GL}}
\def\Stab{\text{Stab}}
\def\p{\mathfrak{p}}
\def\b{\mathfrak{b}}
\def\m{\mathfrak{m}}
\def\O{\mathcal{O}}
\def\q{\mathfrak{q}}
\def\G{\mathbb{G}}
\def\Res{\mathop{\mathrm{Res}}\nolimits}
\def\Ext{\mathop{\mathrm{Ext}}\nolimits}
\def\Alg{\mathop{\mathrm{Alg}}\nolimits}
\def\Stab{\mathop{\mathrm{Stab}}\nolimits}
\def\A{\mathcal{A}}
\def\ind{\mathop{\mathrm{ind}}\nolimits}
\def\Irr{\mathop{\mathrm{Irr}}\nolimits}
\def\add{\mathop{\mathrm{add}}\nolimits}
\def\rem{\mathop{\mathrm{rem}}\nolimits}
\def\ch{\mathop{\mathrm{ch}}\nolimits}
\def\Vec{\mathrm{Vec}}
\def\sVec{\mathrm{sVec}}
\def\Ver{\mathrm{Ver}}
\def\k{\mathbf{k}}
\def\ev{\mathrm{ev}}
\def\coev{\mathrm{coev}}
\def\Rep{\mathop{\mathrm{Rep}}\nolimits}
\def\inv{\mathrm{inv}}
\def\coker{\mathop{\mathrm{coker}}\nolimits}
\def\length{\mathop{\mathrm{length}}\nolimits}
\def\gr{\mathop{\mathrm{gr}}\nolimits}
\def\SFPdim{\mathop{\mathrm{SFPdim}}\nolimits}
\def\FPdim{\mathop{\mathrm{FPdim}}\nolimits}
\def\sdim{\mathop{\mathrm{sdim}}\nolimits}
\def\Gal{\mathop{Gal}}


\pagestyle{plain}

\begin{document}

\title{Twisted Deligne products of semisimple tensor categories}
\author{Pavel Etingof}
\address{Department of Mathematics
Massachusetts Institute of Technology
\newline
77 Massachusetts Avenue,
Cambridge, MA 02139,
USA
}
\email{etingof@math.mit.edu}


\author{Dmitri Nikshych}
\address{Department of Mathematics
University of New Hampshire
\newline Durham, NH 03824}
\email{ Dmitri.Nikshych@unh.edu}

\author{Victor Ostrik}
\address{Department of Mathematics, 
University of Oregon
\newline
Eugene, OR 97403
}
\email{vostrik@uoregon.edu} 

\begin{abstract}
 We discuss the classification of twisted Deligne products of two semisimple tensor categories $\C,\D$, i.e., categorifications of the tensor product of their Grothendieck rings in which the factors
are categorified by $\C$ and $\D$. In particular, we show that if  both factors have no non-trivial gradings, or if one factor has neither non-trivial gradings nor tensor structures on the identity functor, then the only twisted Deligne product is the ordinary one. Using the work \cite{MPP}, this gives, in principle, a group-theoretical classification of twisted Deligne products and, more generally, exact factorizations of arbitrary fusion categories. In the Appendix we introduce the notion of categorical $n$-cocycles for $n=2,3,4$ and show that they are all pullbacks of group $n$-cocycles from the universal grading group of the underlying based ring. In the case of $4$-cocycles, this answers a question of Johnson-Freyd, Ostrik and Yu. 
\end{abstract} 

\maketitle

\centerline{\bf To George Lusztig for his 80th birthday with admiration}

\tableofcontents

\section{Introduction} 

Let $\bold k$ be an algebraically closed field. 
Let $\C,\D,\E$ be semisimple tensor categories over $\bold k$ (\cite{EGNO}, Definition 4.1.1).
Following \cite{G}, 
we say that $\E=\C\bullet \D$ is an {\it exact factorization} if $\E$ is equipped with tensor subcategory embeddings $\C,\D\hookrightarrow \E$ such that the tensor product defines an abelian equivalence $\C\boxtimes \D\to \E$. It is shown in \cite{G} that this is equivalent to the conditions that $\C\D=\E$ and $\C\cap \D=\Vec$.

Furthermore, we say that $\E=\C\bullet \D$ is a {\it twisted Deligne product} if 
$$
X\otimes Y\cong Y\otimes X,\ X\in \C,\ Y\in \D.
$$
That is, a twisted Deligne product of $\C,\D$ is a categorification of the tensor product ${\rm Gr}(\C)\otimes {\rm Gr}(\D)$ of the Grothendieck rings of $\C,\D$ in which the subcategories 
corresponding to the factors are identified with $\C,\D$, respectively. Twisted Deligne products of $\C,\D$ 
are distinguished up to equivalence that restricts to identity tensor functors on $\C$ and $\D$. The simplest example of a twisted Deligne product is the usual Deligne product, $\E=\C\boxtimes \D$. 

\begin{example}\label{groupprod} For a group $L$ and $\omega\in H^3(L,\bold k^\times)$, denote by $\Vec_L^\omega$ 
the tensor category of $L$-graded vector spaces with associativity defined by a cocycle representing $\omega$. Let $L=GH$ be an exact factorization of groups, and $\omega\in H^3(L,\bold k^\times)$. 
Let $\omega_G,\omega_H$ be the restrictions of $\omega$ to $G,H$. Then $\Vec_L^\omega=\Vec_G^{\omega_G}\bullet \Vec_H^{\omega_H}$ is an exact factorization. Moreover, this is a twisted Deligne product iff $G,H$ commute, i.e., $L=G\times H$.
\end{example} 
 
The main goal of this paper is to discuss classification of twisted Deligne products. Namely, one of our main results 
is the following theorem. 

For a tensor category $\mathcal C$, let $U_{\mathcal C}$ be its universal grading group 
(\cite[3.2]{GN}, \cite[Subsection 4.14]{EGNO}), and let ${\rm Aut}_1(\C)$ be the group of isomorphism classes of tensor structures on the identity functor ${\rm Id}_\C$.

 \begin{theorem} \label{maint} (i) (Theorem \ref{part1}) If $U_\C= U_\D=1$ then 
 any twisted Deligne product $\C\bullet \D$ is equivalent to the usual Deligne product $\C\boxtimes \D$ 
 (tautologically on $\C$ and $\D$). 
 
 (ii) (Theorem \ref{part2}) 
 The same holds if $U_\C=1$ and $\Aut_1(\C)=1$ or $U_\D=1$ and $\Aut_1(\D)=1$.
 \end{theorem} 
 
 As indicated in \cite{MPP}, Remark 3.27, a classification of twisted Deligne products of fusion categories in principle allows one to classify all their exact factorizations in group-theoretical terms using the extension theory of \cite{ENO}. In fact, for this one doesn't even need a full classification of twisted Deligne products: Theorem \ref{maint}(i) is already sufficient. Namely, let $\C_{\rm ad}$ denote the adjoint subcategory of $\C$ (\cite{EGNO}, Subsection 4.14). In Corollary 3.24 and Proposition 3.25 of \cite{MPP} it is shown that 
if $\E=\C\bullet \D$ is an exact factorization of fusion categories then $\E_{\rm ad}=\C_{\rm ad}\bullet \D_{\rm ad}$ is a twisted Deligne product. The category $\E$ can then be recovered from $\E_{\rm ad}$ as an extension of $\E_{\rm ad}$ by the universal grading group $U_\mathcal E$, and such extensions can in principle be classified as in \cite{ENO} if one knows the Brauer-Picard groupoid of $\E_{\rm ad}$. 

Now, we can replace $\C,\D$ with $\C_{\rm ad},\D_{\rm ad}$ and iterate this process until we reach a situation when $\C=\C_{\rm ad}$, $\D=\D_{\rm ad}$, i.e., $U_\C=U_\D=1$ (this must eventually happen because the Frobenius-Perron dimension of $\E$ decreases by an integer factor $|U_\E|>1$ at each step).
 But in this case the twisted Deligne product is the ordinary one by Theorem \ref{maint}(i). Thus this provides a recursive group-theoretical classification of exact factorizations of fusion categories, in particular of their twisted Deligne products.

Admittedly, this classification is a bit implicit. After all, computation of Brauer-Picard groups of Deligne products is a rather nontrivial problem, as is classifying group extensions of a given category, and the answer can be complicated. However, in some cases the answer is very simple. For example, Theorem \ref{maint}(i) 
and the results of \cite{MPP} imply that if $U_\C=U_\D=1$ then every exact factorization 
$\C\bullet \D$ is the ordinary Deligne product $\C\boxtimes \D$. 

The paper is organized as follows. Section 2 contains preliminaries. In Section 3 we give a partial classification of twisted Deligne products. Finally, in the Appendix we introduce the notion of categorical $n$-cocycles for $n=2,3,4$ and show that they are all pullbacks of group $n$-cocycles on the universal grading group of the underlying based ring. In the case of $4$-cocycles, this answers a question of Johnson-Freyd, Ostrik and Yu. 

{\bf Acknowledgements.} In this paper, especially in the Appendix, we collaborated with ChatGPT 5.5 Pro. The work of P.E. was supported by the NSF grants DMS-2001318 and DMS-2502467. 
The work of D.N.\ was supported  by  the  National  Science  Foundation  under  Grant No.\ DMS-2302267.

\section{Preliminaries} 

\subsection{The universal grading group} Let $R$ be a based ring with basis $I$  (\cite{EGNO}, Subsection 3.1). Then $R$ has partial order: $a\le b$ if $b-a$ has non-negative coefficients.

Introduce a relation on $I$ by setting $x\sim y$ if there exist $u_1,...,u_n\in I$ such that $x,y\le u_1...u_n$. 
This is an equivalence relation: if $x\sim y$ and $y\sim z$ 
then $x,y\le u_1...u_n$ and $y,z\le v_1...v_m$, hence
$$
x\le xy^*y\le u_1...u_ny^*v_1...v_m,\ z\le yy^*z\le u_1...u_ny^*v_1...v_m.
$$
Let $U(R):=I/\sim$ and denote the equivalence class of $x\in I$ by $\deg x$. 
Thus $I=\sqcup_{g\in U(R)}I_g$, where $I_g$ is the corresponding equivalence class in $I$. It is clear that if $x,y,z\in I$ 
and $x\sim y$, then for any $v,w\in I$ with $v\le xz$, $w\le yz$ or $v\le zx$, $w\le zy$, we have $v\sim w$.
Thus we have a well-defined associative operation on $U(R)$ given by $\deg x\cdot \deg y=\deg z$ if $z\le xy$. Moreover, $1=\deg\bold 1$ is a unit for this operation. Finally, 
if $x\sim y$ then $x^*\sim y^*$, so for all $g\in U(R)$ there is a well-defined 
inverse $g^{-1}=g^*$. Thus $U(R)$ is a group. It is called the {\it universal grading group} of $R$
 (\cite[3.2]{GN}, \cite[Subsection 4.14]{EGNO}). Thus the based ring $R$ has an $U(R)$-grading: $R=\oplus_{g\in U(R)}R_g$, where $R_g:=\Bbb ZI_g$. 
 
The group $U(R)$ has the following universal property. 

\begin{lemma}\label{homom}  Let $G$ be a group and $f : I\to G$ a function 
such that $f(y)=f(x)f(z)$ whenever $y\le xz$. Then there is a unique homomorphism $\overline{f}: U(R)\to G$ such that $f(x)=\overline{f}(\deg x)$.
\end{lemma}

\begin{proof} Let $x,y\le u_1\cdots u_n$. We prove that $f(x)=f(y)$ by induction in $n$. 
The base $n=1$ is trivial. Let $n>1$. Then there exist $x',y'\le u_1...u_{n-1}$ such that $x\le x'u_n,y\le y'u_n$. 
Thus $f(x')=f(y')$ by the induction assumption, so by the hypothesis $f(x)=f(x')f(u_n)=f(y')f(u_n)=f(y)$. Hence $f$ is constant on $\sim$-classes and descends uniquely to $U(R)$.
\end{proof}

In particular, faithful $G$-gradings on $R$ correspond to surjective homomorphisms
$U(R)\to G$. 

If $\C$ is a semisimple tensor category and $R={\rm Gr}(\C)$ is the Grothendieck ring of $\C$, 
then we will call $U(R)$ the {\it universal grading group of $\C$} and denote it by $U_\C$. 
In this case we have $\C=\oplus_{g\in U_\C}\C_g$, where $\C_g$ is spanned by simple objects $X$ with 
$\deg X=g$. Thus faithful $G$-gradings on $\C$ correspond to surjective homomorphisms 
$U_\C\to G$. 

\subsection{The group ${\rm Aut}_1(\C)$}
Let $\C$ be a semisimple tensor category over $\bold k$ with the set of simple objects ${\rm Irr}(\C)$. Let 
$\Aut(\C)$ be the group of tensor autoequivalences of $\C$ up to isomorphism. For a tensor autoequivalence $\phi:\C\to \C$, denote its class in $\Aut(\C)$ by $\widehat\phi$. Let $\Aut_1(\C)\subset \Aut(\C)$ be the subgroup of autoequivalences acting trivially on ${\rm Gr}(\C)$, i.e., 
tensor structures on the identity functor of $\C$ (after choosing isomorphisms $X\cong \phi(X)$ for simple $X\in \C$).

\begin{remark}\label{softaut}
The group $\Aut_1(\C)$ has appeared in the literature under several names. It was studied categorically by Davydov \cite{D1,D2} as the group of twisted forms, or tensor deformations, of the identity functor. In the Hopf algebra setting, the corresponding cohomology group was introduced by Schauenburg \cite{S}; it was subsequently called the second \emph{lazy cohomology} group and systematically studied by Bichon and Carnovale \cite{BC}. More precisely, for a finite-dimensional Hopf algebra $H$, the group $\Aut_1(\Rep(H))$ is identified with the group of invariant twists on $H$ modulo gauge equivalence, or, equivalently, with the lazy cohomology group of $H^*$. Davydov later called its elements \emph{soft} tensor autoequivalences \cite{D3}. Guillot and Kassel \cite{GK} showed that this group need not be abelian and gave such examples already for $\C=\Rep(G)$, where $G$ is a finite group (\cite{GK}, Proposition 1.4).
\end{remark}

We have a natural homomorphism
$$
\theta: H^2(U_\C,\bold k^\times)\to \Aut_1(\C)
$$ 
given by $\theta([\beta])|_{W_1\otimes W_2}=\beta(\deg W_1,\deg W_2)$, where $\beta$ is a normalized 2-cocycle on $U_\C$.  

\begin{lemma}\label{inje}  $\theta$ is injective. Thus it identifies $H^2(U_\C,\bold k^\times)$ 
with a central subgroup of $\Aut_1(\C)$.
\end{lemma}

\begin{proof} If $\theta([\beta])=1$ then $\beta(\deg W_1,\deg W_2)=f(W_1)f(W_2)f(W_3)^{-1}$ for each $W_3\subset W_1\otimes W_2$, where $f: {\rm Irr}(\C)\to \bold k^\times$. Thus, if $X$ occurs in $W_1\otimes\cdots\otimes W_n$ then 
$f(X)^{-1}\prod_i f(W_i)$ depends only on the degrees $\deg W_i$ of the tensor factors. So if $X,X'\subset W_1\otimes\cdots\otimes W_n$ then $f(X)=f(X')$. Hence $f(X)=\overline{f}(\deg X)$, where $\overline{f}$ is a function on $U_\C$, and $\beta=d\overline{f}$. Also, the image of $\theta$ acts on each simple summand of $X\otimes Y$ by the scalar $\beta(\deg X,\deg Y)$, so it commutes with any tensor structure on the identity functor, implying that ${\rm Im}\theta$ is central.
\end{proof} 

However, $\theta$ need not be surjective since according to Remark \ref{softaut}, $\Aut_1(\C)$ can be non-abelian. 

\subsection{Twisted Drinfeld center} 

Recall that the {\it Drinfeld center} of a tensor category $\C$ may be defined as 
the category of $\C$-bimodule functors $\C\to \C$. More generally, 
let $\phi: \C\to \C$ be a tensor autoequivalence. Then 
we can define the $\C$-bimodule category $\C(\phi)$ 
which is $\C$ as a left $\C$-module category, while the right $\C$-action is twisted by $\phi$. That is,  
for $Y\in \C(\phi)$ and $X\in \C$, $Y\circ X:=Y\otimes \phi(X)$. 
Now define the {\it twisted Drinfeld center} of $\C$ to be the
category of $\C$-bimodule functors $\C\to \C(\phi)$. Then, if $\psi: \C\to \C$ 
is another tensor autoequivalence, then the
category of $\C$-bimodule functors $\C(\psi)\to \C(\phi)$ is $\Z_{\phi\psi^{-1}}(\C)$. 

Recall that any module endofunctor of $\C$ as a left $\C$-module is right multiplication 
by an object $Z$ of $\C$. Thus in concrete terms, an object 
of $\Z_\phi(\C)$ can be realized as an object $Z$ of $\C$ together with a functorial isomorphism 
$c_X: X\otimes Z\to Z\otimes \phi(X)$, $X\in \C$, satisfying the consistency relation
\begin{equation}\label{conrel}
c_{X\otimes Y}=(\id_Z\otimes \mathcal J_{X,Y})(c_X\otimes\id_{\phi(Y)})(\id_X\otimes c_Y),
\end{equation}
where $\mathcal J_{X,Y}:\phi(X)\otimes\phi(Y)\to \phi(X\otimes Y)$ is the tensor structure of $\phi$ and we suppress associativity isomorphisms (in particular, $c_{\bold 1}=\id_Z$). Namely, a functor $F\in \Z_\phi(\C)$ gives rise to the object $Z_F:= F(\bold 1)$, and conversely, $F_Z(X)=X\otimes Z$.

Let $G$ be a group and $\beta: G\times G\to \k^\times$ be a 2-cocycle. Denote by 
$\Rep_\beta(G)$ the category of finite dimensional projective representations $V$
of $G$ with $2$-cocycle $\beta$, i.e. maps $\rho: G\to GL(V)$ 
such that 
$$
\rho(gh)=\rho(g)\rho(h)\beta(g,h),\ g,h\in G.
$$ 

\begin{lemma}\label{twistedcenter} Let $\phi:\C\to \C$ be a tensor autoequivalence and let \linebreak $\theta: H^2(U_\C,\k^\times)\to {\rm Aut}_1(\C)$ be the map of Lemma \ref{inje}. 
Then there exists nonzero $Z\in \Z_\phi(\C)$ such that the underlying object of $Z$ is $V\otimes \bold 1$ for a finite dimensional $\k$-vector space $V$ if and only if 
$$
\widehat{\phi}=\theta([\beta])\in \Aut_1(\C)\subset \Aut(\C)
$$ 
for some 2-cocycle $\beta$ on $U_\C$ with values in $\k^\times$ for which $\Rep_\beta(U_\C)\ne 0$. Moreover, $\beta$ is the same for all $Z$ up to coboundaries and can be chosen so that $\beta^{\dim V}=1$. Finally, in this case the category of such $Z$ is naturally equivalent to 
$\Rep_\beta(U_\C)$.
\end{lemma}

\begin{proof} For a simple $X\in \C$ the isomorphism 
$c_X:X\otimes Z\to Z\otimes\phi(X)$ 
gives an isomorphism $V\otimes X\cong V\otimes\phi(X)$. Hence $\phi(X)\cong X$.  Choose such identifications for all simple $X$ and let $\mathcal J$ denote the resulting tensor structure on $\Id_\C$.
Then $c_X$ becomes an operator $\zeta_X\in GL(V)$ for every simple $X$.
For a simple summand $W_3\subset W_1\otimes W_2$, \eqref{conrel} gives
\begin{equation}\label{concon}
\zeta_{W_3}\otimes \id=\zeta_{W_1}\zeta_{W_2}\otimes \gamma(W_1,W_2;W_3)
\end{equation}
on $V\otimes \Hom(W_3,W_1\otimes W_2)$, where $\gamma(W_1,W_2;W_3)$ is the operator by which $\mathcal J_{W_1,W_2}$ acts on $\Hom(W_3,W_1\otimes W_2)$.  It follows that $\gamma(W_1,W_2;W_3)$ is a scalar.

For $g\in GL(V)$, let $[g]$ be the image of $g$ in $PGL(V)$. Since $\gamma$ is a scalar, we have 
\[
[\zeta_{W_3}]=[\zeta_{W_1}][\zeta_{W_2}]
\]
whenever $W_3\subset W_1\otimes W_2$. By Lemma~\ref{homom}, it follows that the assignment $X\mapsto[\zeta_X]$ factors through a homomorphism $U_\C\to PGL(V)$.  For $g\in U_\C$ choose lifts $\rho(g)\in SL(V)$ of this projective representation with $\rho(1)=\id_V$ (which is possible since $\k$ is algebraically closed) and define a 2-cocycle $\beta$ on $U_\C$ by
\[
\rho(gh)=\rho(g)\rho(h)\beta(g,h),\ g,h\in U_\C.
\]
Taking determinants gives $\beta(g,h)^{\dim V}=1$ for all $g,h$.

After rescaling the chosen identifications $\phi(X)\cong X$ by some scalars $\lambda_X\in \k^\times$, we may assume that $\zeta_X=\rho(\deg X)$.
Then \eqref{concon} implies
\[
\gamma(W_1,W_2;W_3)=\beta(\deg W_1,\deg W_2)
\]
whenever $W_3\subset W_1\otimes W_2$.  Hence the tensor structure of $\phi$ is exactly the image of $[\beta]$ under $\theta$, i.e. $\widehat\phi=\theta([\beta])$. In particular, if $Z,Z'\in \mathcal Z_\phi(\C)$ 
are nonzero and give rise to cocycles $\beta,\beta'$ then $[\beta]=[\beta']$. 

Conversely, suppose that, after identifying $\phi$ with the identity on objects, $\mathcal J_{W_1,W_2}$ acts on $W_1\otimes W_2$ by $\beta(\deg W_1,\deg W_2)$.  If $V$ is a projective representation of $U_\C$ with 2-cocycle $\beta$ (i.e. the action $\rho$ of $U_\C$ on $V$ satisfies $\rho(gh)=\rho(g)\rho(h)\beta(g,h)$), set $Z=V\otimes\bold 1$ and define $c_X$ as $\rho(\deg X)$ (upon the chosen identification $\phi(X)\cong X$).  Then the cocycle identity for $\beta$ and the formula for $\mathcal J$ give precisely the consistency relation \eqref{conrel}.  Morphisms between such objects are exactly linear maps intertwining all $\rho(g)$, so the resulting subcategory of $\Z_\phi(\C)$ is equivalent to the category of projective representations of $U_\C$ with 2-cocycle $\beta$.
\end{proof}

\subsection{The K\"unneth theorem}

Let $G,H$ be groups and let $A$ be an abelian group, regarded as a trivial $G\times H$-module.  For $n>0$, define restriction maps 
$$
\xi_{0,n}: H^n(G\times H,A)\to H^n(G,A),\ \xi_{n,0}: H^n(G\times H,A)\to H^n(H,A).
$$ 
Let 
\[
\Gamma_n(G,H,A):=\ker(\xi_{0,n}\oplus \xi_{n,0}).
\]
The two projections $G\times H\to G,H$ give canonical sections of $\xi_{n,0},\xi_{0,n}$,
 so for $n>0$ one has a canonical decomposition
\[
H^n(G\times H,A)\cong H^n(G,A)\oplus H^n(H,A)\oplus \Gamma_n(G,H,A).
\]

The split extension
\[
1\to G\to G\times H\to H\to 1
\]
gives a spectral sequence with the $E_2$-page
\[
E_2^{p,q}=H^p(H,H^q(G,A))\Rightarrow H^{p+q}(G\times H,A),
\]
which is a special case of the Lyndon--Hochschild--Serre spectral sequence for the cohomology of an extension
$L$ of $H$ by $G$. Namely, it is the spectral sequence of a filtered cochain complex computing $H^*(L,A)$, and its $E_\infty$-page is the associated graded object of the induced filtration on cohomology. 

In our case, when $L=G\times H$ is a direct product, the Lyndon--Hochschild--Serre spectral sequence degenerates at $E_2$.
 Thus, 
after removing the two restriction summands, we obtain a natural filtration on $\Gamma_n(G,H,A)$ with associated graded pieces
\[
H^p(H,H^{n-p}(G,A)),\  1\le p\le n-1.
\]
Equivalently, there are natural surjective maps, defined successively on kernels,
\[
\xi_{1,n-1}:\Gamma_n(G,H,A)\to H^1(H,H^{n-1}(G,A)),
\]
\[
\xi_{2,n-2}:\ker\xi_{1,n-1}\to H^2(H,H^{n-2}(G,A)),
\]
and so on. Namely, the Alexander--Whitney and Eilenberg--Zilber comparison maps identify the filtered bar complex with the tensor product of the bar complexes for $G$ and $H$, so that this filtration is identified with the standard K\"unneth filtration, and identify the maps $\xi_{i,n-i}$ with the corresponding edge maps. See \cite[Chapter VII, Sections 6--7]{Brown} for more details.

\begin{example}\label{ex1} Recall that if a group $L$ acts trivially on a group $M$ then $H^1(L,M)=\Hom(L,M)$. 
Therefore the groups $\Gamma_2$ and $\Gamma_3$ have the following structure. 

1. One has \scriptsize
$$
\Gamma_2(G,H,A)=\Hom(H,\Hom(G,A))=\Hom(G,\Hom(H,A))=\Hom(G_{\rm ab}\otimes H_{\rm ab},A),
$$ \normalsize
where the subscript denotes abelianization. 

2. There are natural short exact sequences
\[
0\to H^2(H,\Hom(G,A))\to \Gamma_3(G,H,A)\xrightarrow{\tau_H} \Hom(H,H^2(G,A))\to 0,
\]
and, after interchanging $G$ and $H$,
\[
0\to H^2(G,\Hom(H,A))\to \Gamma_3(G,H,A)\xrightarrow{\tau_G} \Hom(G,H^2(H,A))\to 0.
\]
The kernel of
\[
\tau_G\oplus\tau_H:\Gamma_3(G,H,A)\to \Hom(G,H^2(H,A))\oplus \Hom(H,H^2(G,A))
\]
may be described as the kernel of the map
\[
\eta: H^2(G,\Hom(H,A))\to \Hom(H,H^2(G,A))
\]
obtained by restricting $\tau_H$ to $\ker\tau_G=H^2(G,\Hom(H,A))$.  On cocycles, this map is represented by transposing the variables:
\[
\eta(\gamma)(h)(g_1,g_2)=\gamma(g_1,g_2)(h).
\]
\end{example}

\section{Towards classification of twisted Deligne products} 

\subsection{Data attached to a twisted Deligne product} 

Let $\E=\C\bullet \D$ be a twisted Deligne product. Then every simple object $Y$ of $\D$ defines a $\C$-bimodule category $\C_Y:=\C\otimes Y$, which is canonically equivalent to $\C$ as a left module via $X\mapsto X\otimes Y$, $X\in \C$. Thus, the right action of $\C$ on $\C_Y$ is given by a tensor autoequivalence $\phi_Y$ of $\C$, i.e., $\C_Y\cong \C(\phi_Y)$, the category $\C$ with the usual left action of $\C$ and the right action twisted by $\phi_Y$. Moreover, $\phi_Y$ acts trivially on the objects of $\C$, so it defines an element of $\Aut_1(\C)$, denoted by $\widehat{\phi}_Y$.

In a similar way, for $X\in \C$ we have an autoequivalence $\psi_X: \D\to \D$ 
which gives rise to an element $\widehat{\psi}_X\in \Aut_1(\D)$. 

Let us fix isomorphisms $\xi_{YX}: Y\otimes X\cong X\otimes Y$ for each simple $X\in \C$ and $Y\in \D$ and extend them additively (hence functorially) to arbitrary objects. These isomorphisms 
define an identification of additive functors $\psi_X\cong {\rm Id}_\D$ and $\phi_Y\cong {\rm Id}_\C$ for all $X,Y$. So we may assume that 
$\psi_X={\rm Id}_\D$ and $\phi_Y={\rm Id}_\C$ as additive functors.

Let $\mathcal J(\phi_Y),\mathcal J(\psi_X)$ be the tensor structures on the identity functors representing $\phi_Y,\psi_X$. 
Then, suppressing the associativity morphisms, we have
\begin{equation} \label{xiiden}
\begin{aligned} 
\xi_{Y,W_1\otimes W_2}=\mathcal J(\phi_Y)_{W_1W_2}\xi_{YW_2}\xi_{YW_1},\\
\xi_{V_1\otimes V_2,X}=\xi_{V_1X}\xi_{V_2X}\mathcal J(\psi_X)_{V_1V_2}.
\end{aligned}
\end{equation}
Indeed, the first identity in \eqref{xiiden} is the commutativity of
\[
\begin{tikzcd}[column sep=large,row sep=large]
Y\otimes (W_1\otimes W_2) \arrow[rr,"{\xi_{Y,W_1\otimes W_2}}"] \arrow[d,"{\alpha^{-1}}"'] && (W_1\otimes W_2)\otimes Y \\
(Y\otimes W_1)\otimes W_2 \arrow[r,"{\xi_{YW_1}\otimes\id}"'] & (W_1\otimes Y)\otimes W_2 \arrow[d,"{\alpha}"] & \\
& W_1\otimes (Y\otimes W_2) \arrow[r,"{\id\otimes\xi_{YW_2}}"'] & W_1\otimes (W_2\otimes Y) \arrow[uu,"{(\mathcal J(\phi_Y)_{W_1,W_2}\otimes\id_Y)\alpha^{-1}}"']
\end{tikzcd}
\]
and the second identity is the commutativity of
\[
\begin{tikzcd}[column sep=large,row sep=large]
(V_1\otimes V_2)\otimes X \arrow[rr,"{\xi_{V_1\otimes V_2,X}}"] \arrow[d,"{\mathcal J(\psi_X)_{V_1,V_2}\otimes\id_X}"'] && X\otimes (V_1\otimes V_2) \\
(V_1\otimes V_2)\otimes X \arrow[r,"{\alpha}"'] & V_1\otimes (V_2\otimes X) \arrow[d,"{\id\otimes\xi_{V_2X}}"] & \\
& V_1\otimes (X\otimes V_2) \arrow[d,"{\alpha^{-1}}"] & \\
& (V_1\otimes X)\otimes V_2 \arrow[r,"{\xi_{V_1X}\otimes\id}"'] & (X\otimes V_1)\otimes V_2 \arrow[uuu,"{\alpha}"'].
\end{tikzcd}
\]

This shows that if we renormalize $\xi_{YX}$ by scalars $\lambda_{YX}\in \bold k^\times$ 
then $\mathcal J(\phi_Y)$ and $\mathcal J(\psi_X)$ change as follows: 
$\mathcal J(\phi_Y)\mapsto \mathcal J'(\phi_Y)$ with 
\begin{equation}\label{tran1} 
\mathcal J'(\phi_Y)_{\Hom(W_3,W_1\otimes W_2)}=\mathcal J(\phi_Y)_{\Hom(W_3,W_1\otimes W_2)}\lambda_{YW_3}\lambda_{YW_1}^{-1}\lambda_{YW_2}^{-1}
\end{equation}
and 
$\mathcal J(\psi_X)\mapsto \mathcal J'(\psi_X)$ with 
\begin{equation}\label{tran2} 
\mathcal J'(\psi_X)_{\Hom(V_3,V_1\otimes V_2)}=\mathcal J(\psi_X)_{\Hom(V_3,V_1\otimes V_2)}\lambda_{V_3X}\lambda_{V_1X}^{-1}\lambda_{V_2X}^{-1}.
\end{equation}

The following proposition gives the universal property of the twisted Deligne product $\E=\C\bullet \D$. 

\begin{proposition}\label{tffromdp} Let $\B$ be a tensor category. Then a tensor functor $F: \E\to \B$ is determined by the following data: 

(1) A tensor functor $F_1: \C\to \B$; 

(2) A tensor functor $F_2: \D\to \B$; 

(3) Isomorphisms 
$$
b_{YX}: F_2(Y)\otimes F_1(X)\cong F_1(X)\otimes F_2(Y),\ X\in {\rm Irr}(\C),Y\in {\rm Irr}(\D),
$$ 
 such that after extending them additively (hence functorially) to all objects and suppressing the associativity and tensor structure morphisms on $F_1,F_2$, we have  
$$
b_{Y,W_1\otimes W_2}=\mathcal J(\phi_Y)_{W_1W_2}b_{YW_2}b_{YW_1},
$$
$$
b_{V_1\otimes V_2,X}=b_{V_1X}b_{V_2X}\mathcal J(\psi_X)_{V_1V_2}.
$$
\end{proposition} 

\begin{proof} Clearly, $F$ determines $F_1$ and $F_2$, and we have 
a canonical identification $F(X\otimes Y)\cong F_1(X)\otimes F_2(Y)$ for simple $X\in \C,Y\in \D$. 
Also, we may identify $F_2(Y)\otimes F_1(X)=F(Y)\otimes F(X)\cong F(Y\otimes X)$ with
$F(X\otimes Y)$ by the isomorphism $F(\xi_{YX})$. Thus, we obtain an isomorphism 
$b_{YX}: F_2(Y)\otimes F_1(X)\cong F_1(X)\otimes F_2(Y)$. 
It is easy to check that it satisfies the identities of (3). 

Conversely, given the data (1)-(3), we can define $F$ by $F(X\otimes Y):=F_1(X)\otimes F_2(Y)$ for simple $X\in \C$, $Y\in \D$ 
and define the tensor structure on $F$ by $\mathcal J_{W_1\otimes V_1,W_2\otimes V_2}=b_{V_1W_2}$ 
(again suppressing the associativity morphisms and tensor structure morphisms on $F_1,F_2$).
Here we identify $V_1\otimes W_2$ with $W_2\otimes V_1$ using the isomorphism $\xi_{V_1W_2}$.  
Then it is easy to check that conditions (3) ensure that $\mathcal J$ satisfies the tensor structure diagram.
Namely, it is sufficient to check this on the simple objects $X_a\otimes Y_b$. Then the faces involving only $\C$ or only $\D$ are exactly the tensor-functor axioms for $F_1$ and $F_2$.  The face with one $Y$ and two $X$'s is precisely the first identity in (3), and the face with two $Y$'s and one $X$ is precisely the second identity in (3).  The remaining mixed cases reduce to these by functoriality and the associativity in the original product $\E$, equivalently by \eqref{xiiden}. 
This proves the proposition.  
\end{proof}

\begin{corollary} \label{dete} 
The structure of the twisted Deligne product $\E=\C\bullet\D$ is completely determined by the functors $\phi_Y$ and $\psi_X$ (i.e., by $\mathcal J(\phi_Y)$ and $\mathcal J(\psi_X)$) for all simple $X\in \C$ and $Y\in \D$. Moreover, the tensor structures related by the transformation defined by \eqref{tran1},\eqref{tran2} give rise to the same twisted Deligne product (up to a tensor equivalence which is the 
identity on $\C$ and $\D$). 
\end{corollary} 

\begin{proof} This follows from the categorical Yoneda lemma, as by Proposition \ref{tffromdp} the universal property of $\E$ is determined by $\mathcal J(\phi_Y)$ and $\mathcal J(\psi_X)$. More concretely, if we have two twisted Deligne products $\E$ and $\E'$ with the same $\mathcal J(\phi_Y)$ and $\mathcal J(\psi_X)$ then by Proposition \ref{tffromdp}, there is a tensor equivalence $F: \E\to \E'$ restricting to the identity on $\C,\D$.  
\end{proof} 

\subsection{Properties of $\phi_Y$ and $\psi_X$}

Let $Y_p$ be the simple objects of $\D$. Then for each $p,q,m$ we have a $\C$-bimodule functor 
$$
F_{pq}^m: \C_{Y_p}\boxtimes_{\C}\C_{Y_q}\to \C_{Y_m}
$$ 
given by the tensor product in $\E$. Namely, upon natural identification of the source with $\C$ 
it reduces to tensoring with the vector space $\Hom(Y_m,Y_p\otimes Y_q)$. Since
 $\C_Y\cong \C(\phi_Y)$, we have 
 $$
 F_{pq}^m: \C(\phi_p)\boxtimes_\C \C(\phi_q)\to \C(\phi_m),
 $$ 
 where $\phi_p:=\phi_{Y_p}$. 

\begin{proposition}\label{cocyc} (i) If $Y_m$ is contained in $Y_p\otimes Y_q$ then $\widehat{\phi}_m=\widehat{\phi}_p\widehat{\phi}_q$ modulo $\Im \theta$. 
Thus, $\widehat\phi$ descends to a homomorphism 
$$
\overline{\phi}: U_\D\to \Aut_1(\C)/H^2(U_\C,\bold k^\times).
$$  

(ii) If $X_k$ is contained in $X_i\otimes X_j$ then $\widehat{\psi}_k=\widehat{\psi}_j\widehat{\psi}_i$ modulo $\Im \theta$. Thus, $\widehat\psi$ descends to an anti-homomorphism 
$$
\overline{\psi}: U_\C\to \Aut_1(\D)/H^2(U_\D,\bold k^\times).
$$  
\end{proposition} 

\begin{proof} (i) Since 
$\C(\phi_p)\boxtimes_\C \C(\phi_q)\cong \C(\phi_p\phi_q)$, the functor $F_{pq}^m$ may be viewed as a $\C$-bimodule functor $\C(\phi_p\phi_q)\to \C(\phi_m)$, i.e., 
it is an object of the twisted Drinfeld center $\Z_{\phi_m\phi_q^{-1}\phi_p^{-1}}(\C)$ which is a multiple of $\bold 1$ as an object of $\C$. Since by Lemma \ref{inje}, $H^2(U_\C,\bold k^\times)\hookrightarrow \Aut_1(\C)$, (i) follows from Lemmas \ref{homom} and  \ref{twistedcenter}. 


(ii) The proof is analogous, but with reversed order of factors. Namely, the relation \(Y\otimes X\cong X\otimes \psi_X(Y)\) implies that when \(X_k\subset X_i\otimes X_j\), we have a $\D$-bimodule functor $\D(\psi_j\psi_i)\to \D(\psi_k)$. 
\end{proof}


\begin{corollary}\label{multfree} If $\D$ is a multiplicity free category (i.e., we have that 
$[Y_p\otimes Y_q:Y_m]\le 1$ for all $p,q,m$), or, more generally, if the multiplicities in $\D$ are coprime to 
the order of $H^2(U_\C,\bold k^\times)$ (e.g., coprime to the order of $U_\C$ when it is finite) then $\widehat\phi$ is a homomorphism 
$U_\D\to \Aut_1(\C)$. 
\end{corollary} 

\begin{proof} We have seen that if $Y_m\subset Y_p\otimes Y_q$ then $\widehat{\phi}_{Y_m}^{-1}\widehat{\phi}_{Y_p}\widehat{\phi}_{Y_q}\in H^2(U_\C,\bold k^\times)$
is the 2-cocycle of a projective representation of $U_\C$ on the vector space $\Hom(Y_m,Y_p\otimes Y_q)$. 
If this space is 1-dimensional or, more generally, 
has dimension coprime to the order of $H^2(U_\C,\bold k^\times)$ (when it is finite) then this cohomology class must be trivial. Then $\widehat\phi$ is a group homomorphism. 
\end{proof}  

Next, if $G,Q$ are groups, $A\subset Q$ a central subgroup, and 
$$
s: G\to Q/A
$$ 
a homomorphism, 
then let $\delta s \in  H^2(G,A)$ be the class of the 2-cocycle
$$
\beta(g,h):=\widetilde s(gh)\widetilde s(h)^{-1}\widetilde s(g)^{-1}
$$
where $\widetilde s: G\to Q$ is a set-theoretical lift of $s$.
Then we have natural connecting maps
$$
\delta_\D: \Hom(U_\D,\Aut_1(\C)/H^2(U_\C,\bold k^\times))\to H^2(U_\D,H^2(U_\C,\bold k^\times)),
$$ 
\[
\begin{aligned}
\delta_\C:
\Hom(U_\C^{\rm op},\Aut_1(\D)/H^2(U_\D,\bold k^\times))
&\to H^2(U_\C^{\rm op},H^2(U_\D,\bold k^\times))\\
&\cong H^2(U_\C,H^2(U_\D,\bold k^\times)).
\end{aligned}
\]
 
\begin{theorem} \label{vani} 
One has $\delta_\D\overline{\phi}=1,\ \delta_\C\overline{\psi}=1$. Thus 
$\overline{\phi}$ comes from a homomorphism
\[
\widetilde{\phi}: U_\D\to \Aut_1(\C)
\]
and $\overline{\psi}$ comes from an anti-homomorphism
\[
\widetilde{\psi}: U_\C\to \Aut_1(\D),
\]
in general non-uniquely.
\end{theorem} 

\begin{proof}
Recall that $\E$ is $U_\C\times U_\D$-graded. View $\E$ as a
$U_\C\times U_\D$-extension of its trivial component $\E_{\rm ad}$, and
apply the usual two-step obstruction theory of \cite{ENO} (this
construction applies to semisimple tensor categories with possibly infinitely many objects). Let
$O_4\in H^4(U_\C\times U_\D,\bold k^\times)$
be the second obstruction. Since the extension $\E$ exists, this
obstruction class vanishes.

We use the filtration of the Lyndon--Hochschild--Serre spectral sequence $E_2^{p,q}=H^p(U_\D,H^q(U_\C,\bold k^\times))
\Longrightarrow H^{p+q}(U_\C\times U_\D,\bold k^\times)$
for the trivial extension
\[
        1\to U_\C\to U_\C\times U_\D\to U_\D\to 1.
\]
The $(0,4)$-component of $O_4$ vanishes because $\C$ is already a tensor
category. The $(1,3)$-component vanishes because each $\C_Y$ is a
$\C$-bimodule category. Hence the first possible mixed component is the
$(2,2)$-component
\[
        \xi_{2,2}(O_4)\in H^2(U_\D,H^2(U_\C,\bold k^\times)).
\]

We claim that
\[
        \xi_{2,2}(O_4)=\delta_\D\overline{\phi}.
\]
To see this, fix once and for all, for every $t\in U_\D$, a simple object
$Y_t\in\D_t$, with $Y_1=\bold 1$, and put $\phi_t:=\phi_{Y_t}$.
Thus $t\mapsto\widehat{\phi}_t$ is a fixed normalized
set-theoretical lift of $\overline\phi$. For each pair $(g,h)$ choose a
simple summand $Z_{g,h}\subset Y_g\otimes Y_h$, put
$\phi_{g,h}:=\phi_{Z_{g,h}}$, and set
\[
        V_{g,h}:=\Hom(Z_{g,h},Y_g\otimes Y_h).
\]
Projection of the tensor product to this summand gives a
$\C$-bimodule functor
\[
        \C(\phi_g\phi_h)\longrightarrow \C(\phi_{g,h})
\]
whose underlying left $\C$-module functor is tensoring by $V_{g,h}$.
By Lemma~\ref{twistedcenter}, it determines a projective
$U_\C$-representation on $V_{g,h}$; choose a normalized multiplier
$\beta_{g,h}\in Z^2(U_\C,\bold k^\times)$. Then
\[
        \theta([\beta_{g,h}])
        =\widehat{\phi}_{g,h}\widehat{\phi}_h^{-1}
         \widehat{\phi}_g^{-1}.
\]
Since $\deg Z_{g,h}=gh=\deg Y_{gh}$, the elements
$\widehat{\phi}_{g,h}$ and $\widehat{\phi}_{gh}$ have the same image
in $\Aut_1(\C)/H^2(U_\C,\bold k^\times)$. Hence, by the injectivity of
$\theta$, there is a unique class
$[\kappa_{g,h}]\in H^2(U_\C,\bold k^\times)$ such that
\[
        \theta([\kappa_{g,h}])
        =\widehat{\phi}_{gh}\widehat{\phi}_{g,h}^{-1}.
\]
Choose a normalized cocycle representative $\kappa_{g,h}$. Since
${\rm Im}\,\theta$ is central, the cocycle
\[
        c_{g,h}:=\kappa_{g,h}\beta_{g,h}
\]
represents
\[
        \widehat{\phi}_{gh}\widehat{\phi}_h^{-1}
        \widehat{\phi}_g^{-1}.
\]
Consequently, $(g,h)\mapsto[c_{g,h}]$ is precisely the connecting
cocycle for the fixed section $t\mapsto\widehat{\phi}_t$.

It remains to identify this connecting cocycle with the
$(2,2)$-component of $O_4$. We use the standard filtered bar-complex
model for the Lyndon--Hochschild--Serre spectral sequence, or
equivalently the Alexander--Whitney comparison with the tensor product
of the bar complexes of $U_\C$ and $U_\D$; see
\cite[Chapter VII, Sections 6--7]{Brown}. With the conventions of
Subsection~2.4, after the $(0,4)$- and $(1,3)$-components have been
killed, the map $\xi_{2,2}$ sends a normalized representative $\Omega$
of $O_4$ to the $U_\D$-2-cocycle with values in the group of 2-cocycles $Z^2(U_\C,\bold k^\times)$
given by
\[
        (g,h)\longmapsto
        \big((a,b)\longmapsto
        \Omega((a,1),(b,1),(1,g),(1,h))\big).
\]

Let $a,b\in U_\C$ and choose simple objects
$X_a\in \C_a$, $X_b\in \C_b$. The scalar
$
        \Omega((a,1),(b,1),(1,g),(1,h))
$
is obtained by evaluating the pentagon defect of the chosen extension
data on the four homogeneous objects
$
        X_a, X_b, Y_g, Y_h.
$
Equivalently, it is the ratio of the two routes in the diagram \scriptsize
\[
\begin{tikzcd}[column sep=huge,row sep=large]
((X_a\otimes X_b)\otimes Y_g)\otimes Y_h
\arrow[r,"\alpha"]
\arrow[d,"\alpha"']
&
(X_a\otimes (X_b\otimes Y_g))\otimes Y_h
\arrow[r,"\alpha"]
&
X_a\otimes ((X_b\otimes Y_g)\otimes Y_h)
\arrow[d,"\id\otimes \alpha"]
\\
(X_a\otimes X_b)\otimes (Y_g\otimes Y_h)
\arrow[rr,"\alpha"']
&&
X_a\otimes (X_b\otimes (Y_g\otimes Y_h)).
\end{tikzcd}
\] \normalsize
Restrict this diagram to the block
\[
        \Hom(X_c,X_a\otimes X_b)\otimes V_{g,h},
        \qquad X_c\subset X_a\otimes X_b.
\]
The block calculation with the pair-dependent lift $\phi_{g,h}$ gives
$\beta_{g,h}(a,b)$: it measures the defect between the right
$\C$-action attached to $Z_{g,h}$ and the two successive right actions
attached to $Y_g$ and $Y_h$. The edge map $\xi_{2,2}$, however, is
computed using the fixed section $t\mapsto\widehat{\phi}_t$. Passing
from the pair-dependent lift $\widehat{\phi}_{g,h}$ to the fixed lift
$\widehat{\phi}_{gh}$ contributes the class
$[\kappa_{g,h}]$. Thus, if
\[
        \Omega_{g,h}(a,b):=
        \Omega((a,1),(b,1),(1,g),(1,h)),
\]
then
\[
        [\Omega_{g,h}]
        =[\kappa_{g,h}\beta_{g,h}]
        =[c_{g,h}]
        \quad\text{in }H^2(U_\C,\bold k^\times).
\]
Equivalently, compatible choices of cocycle representatives make the
first equality hold already at the cochain level. Hence
$\xi_{2,2}(O_4)$ is represented by
$(g,h)\mapsto[c_{g,h}]$, and therefore
\[
        \xi_{2,2}(O_4)=\delta_\D\overline\phi.
\]

Since $O_4=1$, we get $\xi_{2,2}(O_4)=1$, and hence
$\delta_\D\overline\phi=1$. The proof of
$\delta_\C\overline\psi=1$ is the same with the two factors interchanged.
\end{proof}


Now suppose $X_k\subset X_i\otimes X_j$ and $Y_m\subset Y_p\otimes Y_q$.  By Proposition~\ref{cocyc} and the description of the central subgroup ${\rm Im}\theta$ in Subsection~2.2, the defect tensor structures
\[
        \phi_m^{-1}\phi_p\phi_q|_{\Hom(X_k,X_i\otimes X_j)}
        \quad\text{and}\quad
        \psi_k^{-1}\psi_j\psi_i|_{\Hom(Y_m,Y_p\otimes Y_q)}
\]
are scalar operators. 

\begin{proposition}\label{phipsi} 
We have 
\begin{equation}\label{equall}
\phi_m^{-1}\phi_p\phi_q|_{\Hom(X_k,X_i\otimes X_j)}=\psi_k^{-1}\psi_j\psi_i|_{\Hom(Y_m,Y_p\otimes Y_q)}.
\end{equation}
\end{proposition} 

\begin{proof} Write $J^\phi_r(i,j;k)$ for the operator by which $\mathcal J(\phi_r)$ acts on the space $\Hom(X_k,X_i\otimes X_j)$ and $J^\psi_s(p,q;m)$ for the operator by which $\mathcal J(\psi_s)$ acts on $\Hom(Y_m,Y_p\otimes Y_q)$. Then the two sides of \eqref{equall} are represented by the scalar defect operators
\[
        J^\phi_m(i,j;k)^{-1}J^\phi_p(i,j;k)J^\phi_q(i,j;k)
\]
and
\[
        J^\psi_k(p,q;m)^{-1}J^\psi_j(p,q;m)J^\psi_i(p,q;m),
\]
respectively.

Choose nonzero morphisms $u:X_k\to X_i\otimes X_j$ and $v:Y_m\to Y_p\otimes Y_q$.  Naturality of $\xi$ gives the commutative square
\[
\begin{tikzcd}[column sep=large,row sep=large]
Y_m\otimes X_k \arrow[r,"{\xi_{Y_mX_k}}"] \arrow[d,"{v\otimes u}"'] & X_k\otimes Y_m \arrow[d,"{u\otimes v}"] \\
(Y_p\otimes Y_q)\otimes (X_i\otimes X_j) \arrow[r,"{\xi_{Y_p\otimes Y_q, X_i\otimes X_j}}"] & (X_i\otimes X_j)\otimes (Y_p\otimes Y_q).
\end{tikzcd}
\]
Let us compare two expansions of the bottom horizontal arrow on the block selected by $u$ and $v$.
First, use the second line of \eqref{xiiden} with $X=X_i$ and with $X=X_j$ to split the crossings of $Y_m$ with $X_i$ and $X_j$ through the summand $Y_m\subset Y_p\otimes Y_q$.  These two applications contribute, in this order, $J^\psi_j(p,q;m), J^\psi_i(p,q;m)$. Then use the first line of \eqref{xiiden} with $Y=Y_m$ and $(W_1,W_2)=(X_i,X_j)$.  On the chosen block this contributes $J^\phi_m(i,j;k)$. 
Thus this expansion gives the operator
\[
        J^\phi_m(i,j;k)\otimes J^\psi_j(p,q;m)J^\psi_i(p,q;m).
\]
Next, use the second line of \eqref{xiiden} with $X=X_k$.  On the chosen block this contributes $J^\psi_k(p,q;m)$.  Then use the first line of \eqref{xiiden} with $Y=Y_p$ and with $Y=Y_q$ to split the crossings of $Y_p$ and $Y_q$ with $X_k$ through the summand $X_k\subset X_i\otimes X_j$.  These two applications contribute, in this order,
$J^\phi_p(i,j;k), J^\phi_q(i,j;k)$. Thus this expansion gives the operator
$$
        J^\phi_p(i,j;k)J^\phi_q(i,j;k)\otimes J^\psi_k(p,q;m)
$$
The commutative square implies equality of these two operators:
\[
        J^\phi_m(i,j;k)\otimes J^\psi_j(p,q;m)J^\psi_i(p,q;m)
        =
        J^\phi_p(i,j;k)J^\phi_q(i,j;k)\otimes J^\psi_k(p,q;m).
\]
Multiplying this equality on the left by $J^\phi_m(i,j;k)^{-1}\otimes\id$ and on the right by $\id\otimes J^\psi_k(p,q;m)^{-1}$ gives \scriptsize
$$
        \id\otimes J^\psi_j(p,q;m)J^\psi_i(p,q;m)J^\psi_k(p,q;m)^{-1}
        =J^\phi_m(i,j;k)^{-1}J^\phi_p(i,j;k)J^\phi_q(i,j;k)\otimes\id.
        $$ \normalsize
Since both Hom spaces under consideration are nonzero, it follows that
\[
J^\psi_j(p,q;m)J^\psi_i(p,q;m)J^\psi_k(p,q;m)^{-1}=\lambda\cdot\id
\]
and
\[
J^\phi_m(i,j;k)^{-1}J^\phi_p(i,j;k)J^\phi_q(i,j;k)=\lambda\cdot\id
\]
for some $\lambda\in \k^\times$.
Now cyclic permutation of factors in the first product gives \eqref{equall}.
\end{proof} 

\subsection{The action of $\Gamma_3(U_\C,U_\D,\k^\times)$ on twisted Deligne products}

\begin{lemma}\label{act0} Let ${\mathcal T}(\C,\D)$ be the set of equivalence classes of twisted Deligne products of $\C$, $\D$. 
Then the group $\Gamma_3(U_\C,U_\D,\k^\times)$ acts naturally on ${\mathcal T}(\C,\D)$. Moreover, this action is free.
\end{lemma} 

\begin{proof} Recall that $\Gamma_3(U_\C,U_\D,\k^\times)$ 
is the kernel of the natural map $H^3(U_\C\times U_\D,\k^\times)\to H^3(U_\C,\k^\times)\oplus H^3(U_\D,\k^\times)$. The action of $\omega\in \Gamma_3(U_\C,U_\D,\k^\times)$ on ${\mathcal T}(\C,\D)$ is by multiplying the associativity morphism on homogeneous objects of degrees $(h_1,g_1),(h_2,g_2),(h_3,g_3)$ by the scalar $\omega((h_1,g_1),(h_2,g_2),(h_3,g_3))$. The restrictions of $\omega$ to $U_\C$ and to $U_\D$ are trivial in cohomology, so after replacing $\omega$ by a cohomologous cocycle if necessary the tensor structures on the embedded subcategories $\C$ and $\D$ are unchanged, so the action is well defined. Since this action modifies $\E\in T(\C,\D)$ as a $U_\E=U_\C\times U_\D$-extension of $\E_{\rm ad}$ by an element of $H^3(U_\E,\k^\times)$, it is free by the classification of extensions in \cite{ENO}. 
\end{proof} 

Now recall from Example \ref{ex1}(2) that we have natural maps \scriptsize
$$
\tau_1: \Gamma_3(U_\C,U_\D,\k^\times)\to H^1(U_\D,H^2(U_\C,\bold k^\times)),\ 
\tau_2: \Gamma_3(U_\C,U_\D,\k^\times)\to H^1(U_\C,H^2(U_\D,\bold k^\times)).
$$
\normalsize
For a (normalized) representative $\omega$, we use the following representatives of these maps.  If $h,h_1,h_2\in U_\C$ and $g,g_1,g_2\in U_\D$, then
$$
\tau_1(\omega)(g)(h_1,h_2)
=
$$
$$
\omega((1,g),(h_1,1),(h_2,1))\omega((h_1,1),(1,g),(h_2,1))^{-1}
 \omega((h_1,1),(h_2,1),(1,g)),
 $$
 $$
\tau_2(\omega)(h)(g_1,g_2)=
$$
$$
\omega((1,g_1),(1,g_2),(h,1))^{-1}\cdot \omega((1,g_1),(h,1),(1,g_2))
 \omega((h,1),(1,g_1),(1,g_2))^{-1}.
$$


\begin{proposition}\label{transfo}
Let $\omega\in \Gamma_3(U_\C,U_\D,\k^\times)$. Then 
the action of $\omega$ from Lemma \ref{act0} transforms $\widehat{\phi}$ and $\widehat{\psi}$ as follows: 
$$
\widehat{\phi}\mapsto \widehat{\phi}\cdot \tau_1(\omega);\ \widehat{\psi}\mapsto \widehat{\psi}\cdot \tau_2(\omega).
$$   
\end{proposition} 

\begin{proof}
The functors $\phi_Y$ and $\psi_X$ are unchanged as additive functors; only their tensor structures change.  For $Y\in\D_g$, compare the first diagram in \eqref{xiiden} before and after multiplying the associator by $\omega$.  The three associators in that diagram contribute respectively
\[
\begin{gathered}
\omega((1,g),(h_1,1),(h_2,1))^{-1},\quad
\omega((h_1,1),(1,g),(h_2,1)),\\
\omega((h_1,1),(h_2,1),(1,g))^{-1}.
\end{gathered}
\]
Therefore $\mathcal J(\phi_Y)$ is multiplied by the inverse of this product, i.e. by $\tau_1(\omega)(g)(h_1,h_2)$.

Similarly, for $X\in\C_h$, the second diagram in \eqref{xiiden} shows that $\mathcal J(\psi_X)$ is multiplied by \scriptsize
\[
\omega((1,g_1),(1,g_2),(h,1))^{-1}
\omega((1,g_1),(h,1),(1,g_2))
\omega((h,1),(1,g_1),(1,g_2))^{-1},
\] \normalsize
which is $\tau_2(\omega)(h)(g_1,g_2)$ with the convention above.  Passing to classes in $\Aut_1$ gives the stated formulas.
\end{proof} 

Note that this induced action on $\widehat\phi,\widehat\psi$ is, in general, not faithful, since the kernel $K:=K(U_\C,U_\D)$ of $\tau_1\oplus \tau_2$ (described in Example \ref{ex1}(2)) acts trivially on these classes.  Thus $\widehat\phi$ and $\widehat\psi$ do not detect the whole $\Gamma_3$-action.  To see the action of $K$, one must keep the actual tensor structures $\mathcal J(\phi_Y)$ and $\mathcal J(\psi_X)$, modulo the renormalizations \eqref{tran1}, \eqref{tran2}.  By Corollary~\ref{dete}, these data determine the twisted product up to equivalence, so it remains to compute how $K$ changes these representatives.

By Example \ref{ex1}(2), the group $K$ may be interpreted as the kernel of the natural map 
$\eta: H^2(U_\D,H^1(U_\C,\bold k^\times))\to H^1(U_\C,H^2(U_\D,\bold k^\times))$ given on cocycles by 
$\eta(\gamma)(h)(g_1,g_2)=\gamma(g_1,g_2)(h)$, $g_1,g_2\in U_\D, h\in U_\C$. 
Thus $K$ consists of images $\gamma_*$ of 2-cocycles $\gamma$ on $U_\D$ with coefficients 
in $\Hom(U_\C,\bold k^\times)$ such that $\gamma(-,-)(h)$ is a coboundary 
for each $h\in U_\C$, but $\gamma$ itself is not necessarily a coboundary. 
Let $\lambda(h,g)$ be a function on $U_\C\times U_\D$ such that 
for each $h$, the cocycle $\gamma(-,-)(h)$ is the coboundary of $\lambda(h,-)$. 
Then $\lambda$ is defined up to multiplication by a character $\mu\in \Hom(U_\D,{\rm Fun}(U_\C,\bold k^\times))$ of $U_\D$ with values in the group of $\k^\times$-valued functions 
on $U_\C$. 

\begin{proposition}\label{act} The action of $\gamma$ on $\phi$ and $\psi$ in Lemma \ref{act0} (up to equivalence) is as follows: 
$\psi$ remains the same and, for $Y\in \D_g$, the tensor structure $\mathcal J(\phi_Y)$ on the block $\Hom(X_3,X_1\otimes X_2)$ is multiplied by the scalar 
$\lambda(h_1,g)\lambda(h_2,g)/\lambda(h_1h_2,g),\ h_a=\deg X_a.$
\end{proposition} 

\begin{proof}
In the normalized bar model we may represent the class of $\gamma$ by the $3$-cocycle
\[
        \omega_\gamma((h_1,g_1),(h_2,g_2),(h_3,g_3))=\gamma(g_1,g_2)(h_3).
\]
With this representative, $\tau_1(\omega_\gamma)=1$ and
$\tau_2(\omega_\gamma)(h)(g_1,g_2)=\gamma(g_1,g_2)(h)^{-1}$.
Choose $\lambda$ so that
\[
        \gamma(g_1,g_2)(h)=\lambda(h,g_1)\lambda(h,g_2)\lambda(h,g_1g_2)^{-1}.
\]
Then the change in $\mathcal J(\psi_X)$ is removed by the renormalization
$$
        \xi_{YX}\mapsto \lambda(\deg X,\deg Y)^{-1}\xi_{YX}.
$$
Using \eqref{tran2}, this renormalization multiplies $\mathcal J(\psi_X)$ by $d\lambda(h,-)$, so $\psi$ returns to its original representative.  By \eqref{tran1}, the same renormalization changes $\mathcal J(\phi_Y)$ by the factor
$
        \lambda(h_1,g)\lambda(h_2,g)/\lambda(h_1h_2,g),
$
as claimed.
\end{proof} 

Note that if we change $\lambda$ to another representative $\lambda'$, then $\lambda/\lambda'=\mu$ is a character with 
respect to $g$. Hence by Corollary \ref{dete}, changing $\phi$ by $\mu(h_1,g)\mu(h_2,g)/\mu(h_1h_2,g)$
is a trivial transformation. 
Also, if $\gamma$ is a coboundary
then $\lambda$ can be chosen a character with respect to $h$, so the action we defined is trivial.  
Thus the action of Proposition \ref{act} is well defined.  

\begin{example} Suppose that $U_{\mathcal C}$ and $U_{\mathcal D}$ are finite and $\k=\mathbb C$.
Put $A=(U_{\mathcal C})_{\rm ab}$ and $B=(U_{\mathcal D})_{\rm ab}$.
The K\"unneth theorem, together with the universal coefficient theorem, identifies the kernel
$K=\ker(\tau_1\oplus \tau_2)$
with
$\operatorname{Hom}(\operatorname{Tor}(A,B),\mathbb C^\times)$.

Equivalently, since the exponential map gives a canonical identification
$\mathbb Q/\mathbb Z \xrightarrow{\sim} \mu_\infty\subset \mathbb C^\times,
\ t\mapsto e^{2\pi i t}$,
and all characters of finite groups have values in $\mu_\infty$, we may write
\[
K\cong \operatorname{Tor}(A,B)^*
\cong A^*\otimes B^*
\cong H^1(U_{\mathcal C},\mathbb C^\times)\otimes H^1(U_{\mathcal D},\mathbb C^\times),
\]
where for a finite abelian group $L$, $L^*:=\operatorname{Hom}(L,\mathbb Q/\mathbb Z)$, which is 
identified with $\Hom(L,\mathbb C^\times)$ by the exponential map. Thus this identification is canonical over $\mathbb C$ (as a topological field).

In terms of this realization, the map $\gamma\mapsto\lambda$ can be described (on decomposable tensors) as follows. 
Let \(c:U_\C\to\mathbb C^\times\) and \(d:U_\D\to\mathbb C^\times\) be characters and $\gamma_*=c\otimes d$.
  Choose functions \(L_c,L_d\), with \(L_c(1)=L_d(1)=0\), such that \(c(h)=\exp L_c(h)\) and \(d(g)=\exp L_d(g)\), and set
\[
        \lambda(h,g)=\exp\!\left(\frac{L_c(h)L_d(g)}{2\pi i}\right).
\]
Then
\[
        \gamma(g_1,g_2)(h)=
        \lambda(h,g_1)\lambda(h,g_2)\lambda(h,g_1g_2)^{-1}
\]
is a representative of the corresponding element of \(K\).  

The dependence on the choices of logarithms is as follows. If
$L_c'(h)=L_c(h)+2\pi iM(h)$, then
\[
        \frac{\lambda'(h,g)}{\lambda(h,g)}=d(g)^{M(h)},
\]
which is a character in the $g$-variable for every fixed $h$ and is
therefore exactly the ambiguity described above. If instead
$L_d'(g)=L_d(g)+2\pi iN(g)$, then
\[
        \nu(h,g):=\frac{\lambda'(h,g)}{\lambda(h,g)}=c(h)^{N(g)}.
\]
For every fixed $g$, the function $\nu(-,g)$ is a character of
$U_\C$, so $d_{U_\C}\nu=1$. Hence the factor
$d_{U_\C}\lambda$ governing the change of $\phi$ is unchanged, while
\[
        \gamma'=d_{U_\D}\lambda'
        =\gamma\,d_{U_\D}\nu
\]
is cohomologous to $\gamma$ as a $U_\D$-cocycle with values in
$H^1(U_\C,\mathbb C^\times)$. Thus the resulting element of $K$, and
the induced action, are independent of the choices of $L_c$ and $L_d$.
\end{example}

\begin{proposition}\label{equii} Let $\E$, $\E'$ be two twisted products with the same $\widehat{\phi}$ and $\widehat{\psi}$. 

(i) If $\C$ or $\D$ is a pointed category then $\E$ and $\E'$ differ by twisting the associativity by an element of $K(U_\C,U_\D)$.

(ii)  If $U_\C$ or $U_\D$ is trivial then $\E\cong \E'$. 
\end{proposition} 

\begin{proof} Let $\phi,\psi$ be the maps corresponding to $\E$ and $\phi',\psi'$ to $\E'$. 
Renormalizing $\xi_{YX}$, we may assume that $\psi=\psi'$. 
Then 
$$
\phi_Y^{-1}\phi_Y'|_{\Hom(W_3,W_1\otimes W_2)}=\lambda_{YW_1}\lambda_{YW_2}\lambda_{YW_3}^{-1}. 
$$
By Proposition \ref{phipsi}, this is multiplicative in $Y$: indeed, applying Proposition~\ref{phipsi} to $\E$ and to $\E'$ and dividing the two equalities cancels the common $\psi$-defect.  Hence by Lemma \ref{homom} 
$$
\lambda_{YW_1}\lambda_{YW_2}\lambda_{YW_3}^{-1}=\zeta(W_1,W_2,W_3)(\deg Y), 
$$
where $\zeta(W_1,W_2,W_3)\in \Hom(U_\D,\bold k^\times)$. 

(i) If $\C$ is pointed, then 
$\zeta$ depends only on $W_1$ and $W_2$ (as $W_3=W_1\otimes W_2$) and is a 2-cocycle on $U_\C$ with coefficients in $\Hom(U_\D,\bold k^\times)$, 
such that for every $g\in U_\D$, the 2-cocycle $\zeta(-,-)(g)$ is a coboundary.
 This implies the statement.

More explicitly, write the simple objects of $\C$ as $X_h$, $h\in U_\C$, and set
\[
        \lambda(h,g):=\lambda_{Y_gX_h},\qquad g\in U_\D.
\]
The displayed equality says that
\[
        \zeta(h_1,h_2)(g)
        =\lambda(h_1,g)\lambda(h_2,g)\lambda(h_1h_2,g)^{-1}.
\]
Since $\zeta(h_1,h_2)$ is a character of $U_\D$, the cochain
\[
        \gamma(g_1,g_2)(h):=
        \lambda(h,g_1)\lambda(h,g_2)\lambda(h,g_1g_2)^{-1}
\]
takes values in $H^1(U_\C,\bold k^\times)$: indeed $d_{U_\C}\gamma=d_{U_\D}\zeta=1$.  Also $d_{U_\D}\gamma=1$, so $\gamma$ is a $2$-cocycle on $U_\D$ with coefficients in $H^1(U_\C,\bold k^\times)$.  For each fixed $h$, the cocycle $\gamma(-,-)(h)$ is the coboundary $d_{U_\D}\lambda(h,-)$, so the image of $[\gamma]$ under the transposition map $\eta$ of Example~\ref{ex1}(2) is $1$.  Hence $[\gamma]\in K(U_\C,U_\D)$.  Proposition~\ref{act}, applied to this representative and to the cochain $\lambda$, leaves $\psi$ unchanged and multiplies $\phi_{Y_g}$ on the $X_{h_1}\otimes X_{h_2}$ block by
\[
        \lambda(h_1,g)\lambda(h_2,g)\lambda(h_1h_2,g)^{-1},
\]
which is exactly the scalar relating the tensor structures of $\E$ and $\E'$.  Therefore the two twisted products differ by the action of the class $[\gamma]\in K(U_\C,U_\D)$.

 If $\D$ is pointed, the proof is symmetrical. 

(ii) If $U_\D$ is trivial then $\zeta=1$ and the statement follows, and if $U_\C$ is trivial then the proof is symmetrical.   
\end{proof} 

\subsection{Examples} 

\begin{example} If $\C=\Vec_G,\D=\Vec_H$, then $U_\C=G$, $U_\D=H$, $\Aut_1(\C)=H^2(G,\bold k^\times)$, $\Aut_1(\D)=H^2(H,\bold k^\times)$. 
Since both categories are multiplicity free, by Corollary \ref{multfree} $\widehat{\phi}$ and $\widehat{\psi}$ are homomorphisms 
$U_\D\to H^2(U_\C,\bold k^\times)$ and $U_\C\to H^2(U_\D,\bold k^\times)$, respectively. 

Also, because of the action of $\Gamma_3(G,H,\k^\times)$, 
any pair of homomorphisms can arise, as $\tau_1\oplus \tau_2$ is surjective.
Indeed since $\bold k$ is algebraically closed, $\k^\times$ is a divisible group; hence
\[
        H^2(G,\k^\times)\cong \Hom(H_2(G,\mathbb Z),\k^\times),
        \qquad
        H^2(H,\k^\times)\cong \Hom(H_2(H,\mathbb Z),\k^\times).
\]
Consider the transposition map from Example~\ref{ex1}(2), 
\[
        \eta:H^2(G,\Hom(H,\k^\times))\to \Hom(H,H^2(G,\k^\times)).
\]
The universal coefficient exact sequence gives a surjection
\[
        H^2(G,\Hom(H,\k^\times))\twoheadrightarrow
        \Hom(H_2(G,\mathbb Z),\Hom(H,\k^\times)).
\]
After currying, the target is canonically
\[
        \Hom(H,\Hom(H_2(G,\mathbb Z),\k^\times))=
        \Hom(H,H^2(G,\k^\times)),
\]
and this is exactly the map induced by $\eta$.  Thus $\eta$ is surjective.  Now choose a class in $\Gamma_3(G,H,\k^\times)$ with the prescribed value of one edge map; the difference between its other edge-map value and the desired one can be corrected by adding an element in the kernel of the first edge map, and the preceding surjectivity of $\eta$ realizes this correction.  Hence $\tau_1\oplus\tau_2$ is surjective.

 Finally, once these homomorphisms 
have been fixed, by Proposition \ref{equii} the possible twisted products are parametrized by a torsor over the group $K$. This torsor is in fact nonempty, since in this case twisted products $\E=\C\bullet \D$ are parametrized by $\Gamma_3(G,H,\k^\times)$ (Example \ref{groupprod}). 

More generally, if $\C=\Vec_G^{\omega_G},\D=\Vec_H^{\omega_H}$, the story is the same, with associator of $\E$ modified by $\omega_G$ and $\omega_H$. This recovers Example \ref{groupprod} in full generality. 
\end{example} 

\begin{example} \label{smashpro}
Let \(\D=\Vec_G\), and denote the simple invertible object corresponding to \(g\in G\) again by \(g\).  If \(\E=\C\bullet \Vec_G\) is a twisted Deligne product, then conjugation by \(g\),
$\phi_g(X)=g\otimes X\otimes g^{-1}, X\in\C$ defines a categorical action \(\phi:G\to\Aut(\C)\) preserving isomorphism classes of simple objects.  

Conversely, any categorical action \(\phi:G\to\Aut(\C)\) preserving simple objects defines a semidirect product \(\C\rtimes \Vec_G\) (\cite{EGNO}, Definition 4.15.5).  Its underlying abelian category is \(\C\boxtimes \Vec_G\),
 and its tensor product is characterized by the crossing relation
$g\otimes X=\phi_g(X)\otimes g$.
This category contains \(\Vec_G\) and \(\C\) as tensor subcategories, and it is a twisted Deligne product of them.  These two constructions are inverse to each other.

Let us compare this with the cohomological invariants above.  Choose, for every simple \(X\in\C\), isomorphisms \(u_{g,X}:\phi_g(X)\to X\).  The monoidal coherence of the action then gives a projective \(G\)-action on \(X\), hence a class $\eta_X\in H^2(G,\k^\times)$
independent of the choices of the \(u_{g,X}\).  Moreover, the induced projective action of \(G\) on \(\Hom(V_3,V_1\otimes V_2)\) has 2-cocycle
$\eta_{V_1}\eta_{V_2}\eta_{V_3}^{-1}$.
Our first invariant of the action is just the homomorphism
$\phi: G\to\Aut_1(\C)$. The other invariant is the collection \(\widehat\psi_X\in H^2(G,\k^\times)\).  The element
$\widetilde\psi_X:=\widehat\psi_X\eta_X^{-1}$ 
depends only on \(\deg X\) and defines a homomorphism
$\widetilde\psi:U_\C\to H^2(G,\k^\times)$ (Theorem \ref{vani}). 
The action of \(\Gamma_3(U_\C,G,\k^\times)\) changes \(\widetilde\psi\) by the edge map \(\tau_2\) (Proposition \ref{transfo}),
so \(\widetilde\psi\) can be prescribed arbitrarily.  For a fixed \(\widehat \phi\), the obstruction to realizing the corresponding object-preserving categorical action lies in \(H^3(G,H^1(U_\C,\k^\times))\); when it vanishes, the choices form a torsor over \(H^2(G,H^1(U_\C,\k^\times))\).  Under the natural map
$H^2(G,H^1(U_\C,\k^\times))\to H^1(U_\C,H^2(G,\k^\times))$
this torsor maps to the possible values of \(\widetilde\psi\), and its kernel is precisely \(K(U_\C,G)\), in agreement with Proposition~\ref{act}.

More generally, we can take $\D=\Vec_G^\alpha$ where $\alpha: G^3\to \k^\times$ is a 3-cocycle. Then the story is the same, with associator of $\E$ modified by $\alpha$.
\end{example}

\begin{remark} Example \ref{smashpro} shows that the action of Proposition \ref{transfo} need not be transitive. Namely, for a finite group $H$ let ${\rm Aut}_c(H)$ be the group 
of automorphisms of $H$ which act trivially on the set of conjugacy classes of $H$, and 
let ${\rm Inn}(H)$ be its normal subgroup of inner automorphisms. As explained in \cite{GK}, 
(proof of Proposition 1.4), there exist $p$-groups $H$ for which  ${\rm Aut}_c(H)/{\rm Inn}(H)$ 
is non-abelian. Let $G:={\rm Aut}_c(H)$, and consider the twisted Deligne product 
$\Rep(H)\rtimes \Vec_G$. Then we have a natural homomorphism $\overline \phi: G\to \Aut_1(\Rep(H))/{\rm Im}\theta$, which is nontrivial since $G/{\rm Inn}(H)\subset \Aut_1(\Rep(H))$ is non-abelian. However, $\overline \phi$ does not change under the action of $\Gamma_3$.
Thus the twisted Deligne product $\Rep(H)\rtimes \Vec_G$ is not equivalent under this action 
to the usual Deligne product $\Rep(H)\boxtimes \Vec_G$.
\end{remark}

\subsection{Classification theorems} 

\begin{theorem}\label{part1} Suppose that $U_\C=U_\D=1$. Then any twisted Deligne product $\E=\C\bullet \D$ is equivalent to the ordinary Deligne product $\C\boxtimes \D$.  
\end{theorem} 

\begin{proof} Since $U_\C=U_\D=1$, the maps $\overline{\phi}$ and $\overline{\psi}$ are trivial. 
Also the groups $H^2(U_\C,\bold k^\times)$ and $H^2(U_\D,\bold k^\times)$ are trivial, so the maps $\widehat{\phi}$ and $\widehat{\psi}$ are trivial. 
By Proposition \ref{equii}(ii), the twisted product $\E$ differs from $\C\boxtimes \D$ by twisting by an element of $K(U_\C,U_\D)=1$. This implies the statement. 
\end{proof} 

We also have the following theorem. 

\begin{theorem}\label{part2} Assume that $U_\C=1$ and $\Aut_1(\C)=1$, or $U_\D=1$ and $\Aut_1(\D)=1$. Then any twisted Deligne product $\E=\C\bullet \D$ is equivalent to the ordinary Deligne product $\C\boxtimes \D$. 
\end{theorem} 

\begin{proof} Let $U_\C=1$ and $\Aut_1(\C)=1$ (the other case is symmetrical). 
Then we have $\widehat\phi=1$. Thus by Proposition \ref{phipsi}, $\widehat \psi_k^{-1}\widehat \psi_j\widehat \psi_i=1$ for $X_k\subset X_i\otimes X_j$.
Hence by Lemma \ref{homom}, $\widehat \psi$ is an antihomomorphism $U_\C\to \Aut_1(\D)$, so $\widehat \psi=1$ since $U_\C=1$. Thus, again using that $U_\C=1$, the result follows from Proposition \ref{equii}(ii). 
\end{proof} 

\section{Appendix: based-ring cocycles and categorical cocycles} 

In this appendix we study based-ring cocycles and categorical cocycles
 in tensor categories, in particular answering Question 4.3 in \cite{JOY}. Although the results below are not used in the main body of the paper, they are motivated 
by some of its arguments, e.g. Lemma \ref{twistedcenter} and its proof. 

\subsection{Based-ring cocycles} 
Let $R$ be a based ring with basis $I$ and structure constants $N_{xz}^y$, and let 
$A$ be an abelian group written multiplicatively. For $n\ge 1$, we call a function
$f:I^n\to A$ a \emph{based-ring $n$-cocycle} if for every $x_1,\ldots,x_{n+1}\in I$ and every choice of $y_i\in I$ with $N^{y_i}_{x_ix_{i+1}}>0$, $1\le i\le n$, one has \scriptsize
\begin{equation}\label{frcocycle}
 f(x_2,\ldots,x_{n+1})
 \prod_{i=1}^{n}
 f(x_1,\ldots,x_{i-1},y_i,x_{i+2},\ldots,x_{n+1})^{(-1)^{i}}
 f(x_1,\ldots,x_{n})^{(-1)^{n+1}}=1 .
\end{equation} \normalsize

\begin{proposition}\label{anydegreecocycle} Let $f:I^n\to A$ be a based-ring $n$-cocycle. Then there is a unique group $n$-cocycle
$\overline f:U(R)^n\to A$
such that
$$
 f(x_1,\ldots,x_n)=\overline f(\deg x_1,\ldots,\deg x_n)
$$
for all $x_1,\ldots,x_n\in I$.
\end{proposition}

\begin{proof} It suffices to show that $f$ is constant on universal grading classes in each variable.  For
$1\le p\le n$ and $m\ge 1$, let $P(p,m)$ be the following assertion: if
$x,x'\le u_1\cdots u_m$, then, for all $1\le q\le p$ and fixed entries
$a_1,\ldots,a_{q-1},a_{q+1},\ldots,a_n$, one has
\begin{equation*}
        f(a_1,\ldots,a_{q-1},x,a_{q+1},\ldots,a_n)
        =
        f(a_1,\ldots,a_{q-1},x',a_{q+1},\ldots,a_n).
\end{equation*}
We shall prove $P(p,m)$ for all $p,m$ by induction on $p$ (starting with the vacuous case $p=0$), and, for each fixed $p$, by induction on $m$.

The case $m=1$ is trivial: $x=x'=u_1$.  So consider $m>1$.  
We assume that $P(q,r)$ has already been proved for all
$q<p$ and all $r$, and that $P(p,r)$ has already been proved for all $r<m$.  

Let
$x,x'\le u_1\cdots u_m$.  Choose elements
$t,t'\le u_1\cdots u_{m-1}$ such that
$
        x\le tu_m,
        \ 
        x'\le t'u_m .
$
Apply \eqref{frcocycle} to the $(n+1)$-tuple
$$
        (a_1,\ldots,a_{p-1},t,u_m,a_{p+1},\ldots,a_n)
$$
with $y_{p}=x$.  Apply it again to
$$
        (a_1,\ldots,a_{p-1},t',u_m,a_{p+1},\ldots,a_n)
$$
with $y_{p}=x'$.  For $i\ne p$, choose the adjacent summands in the two equations as follows.  If the adjacent pair does not involve $t$ or $t'$, choose the same summand in the two equations.  If $i=p-1$ (when $p>1$), choose arbitrary summands $
        y_{p-1}\le a_{p-1}t,
        \ 
        y'_{p-1}\le a_{p-1}t'.
$


We compare the two cocycle equations. The only factors not yet known to agree are the \(i=p\) factors:
$f(a_1,\ldots,a_{p-1},x,a_{p+1},\ldots,a_n)$
in the first equation and
$
        f(a_1,\ldots,a_{p-1},x',a_{p+1},\ldots,a_n)
$
in the second one.  We claim that every other factor in the first equation equals the corresponding factor in the second equation.

Indeed, consider first the initial factor $f(x_2,\ldots,x_{n+1})$.
If $p=1$, this factor is the same in the two equations.  If $p>1$, then $t$ and $t'$ occur in this factor in position $p-1$; since $t,t'\le u_1\cdots u_{m-1}$, equality follows from $P(p-1,m-1)$.

Next consider a middle factor with $i<p-1$ (when $p>2$).  After replacing $x_i,x_{i+1}$ by the chosen summand $y_i$, the symbols $t$ and $t'$ occur in position $p-1$ of this factor.  Thus the two corresponding factors are equal by $P(p-1,m-1)$.  For $i=p-1$ (when $p>1$), the two corresponding factors contain $y_{p-1}$ and $y'_{p-1}$ in position $p-1$.  But both $y_{p-1}$ and $y'_{p-1}$ occur in the product
$a_{p-1}u_1\cdots u_{m-1}$.
Hence equality follows from $P(p-1,m)$.

For every middle factor with $i>p$, the symbols $t$ and $t'$ occur in position $p$ of the factor, while all other entries are the same.  Since
$t,t'\le u_1\cdots u_{m-1}$, equality follows from $P(p,m-1)$.

It remains only to compare the final factor $f(x_1,\ldots,x_n)$.
In this factor the symbols $t$ and $t'$ occur in position $p$, so equality again follows from $P(p,m-1)$.

Thus all factors in the two instances of \eqref{frcocycle} agree except possibly the $i=p$ factors.  Since both products equal $1$, the $i=p$ factors also agree.  Therefore
\begin{equation*}
        f(a_1,\ldots,a_{p-1},x,a_{p+1},\ldots,a_n)
        =
        f(a_1,\ldots,a_{p-1},x',a_{p+1},\ldots,a_n),
\end{equation*}
which proves $P(p,m)$.  This completes the double induction.

It follows that $\overline f$ is a group cocycle. Indeed, for $g_1,\ldots,g_{n+1}\in U(R)$ choose $x_r\in I_{g_r}$ and $y_i\le x_i x_{i+1}$; then $\deg y_i=g_i g_{i+1}$, so \eqref{frcocycle} becomes exactly the group-cocycle identity. 
\end{proof} 

\subsection{Categorical $2$- and $3$-cocycles} 

Let $R$ be a based ring with basis $I$ and structure constants $N_{ij}^k$. Recall that a {\it categorification} of $R$ over $\k$ is a semisimple tensor category $\C$ with an identification ${\rm Gr}(\C)\cong R$; i.e., the simple objects of $\C$ are $X_i$, $i\in I$, with $X_i\otimes X_j\cong \oplus_k N_{ij}^k X_k$ (\cite{EGNO}, Subsection 4.10).  

\begin{definition} Let $\C$ be a categorification of $R$ with associativity isomorphism $\alpha$. 

(i) A function $f:I^2\to \k^\times$ is called a {\it categorical $2$-cocycle} for $\C$ if the morphisms
$f(i,j)\id_{X_iX_j}$ extended additively define a tensor structure on the identity functor of $\C$.

(ii) A function $f:I^3\to \k^\times$ is called a {\it categorical $3$-cocycle} for $\C$ if the morphisms
$f(i,j,k)\alpha_{X_iX_jX_k}$ extended additively define a new associativity constraint on $\C$. 
\end{definition} 

\begin{example}\label{categorical-cocycle-examples} (i) The pullback to $I$ of a group 2-cocycle on $U_\C$ is a categorical 2-cocycle for $\C$. The pullback to $I$ of a group $3$-cocycle on $U_\C$ is a categorical 3-cocycle for $\C$. 

(ii) Let $\D$ be a tensor category and $F: \C\to \D$ a tensor functor with tensor structure $\mathcal J$. Say that a collection of $h_i\in \Aut F(X_i),i\in I$ defines a {\it 
projective tensor automorphism} of $F$ if there exist scalars $\beta(i,j)\in \k^\times,\ i,j\in I$ such that 
$$
\mathcal J_{X_iX_j}(h_i\otimes h_j)\mathcal J_{X_iX_j}^{-1}|_{F(X_k)}=\beta(i,j)h_k
$$ 
for all $i,j\in I$ and $X_k\hookrightarrow X_i\otimes X_j$. It is easy to see that 
if $h$ is a projective tensor automorphism then $\beta$ is a categorical 2-cocycle for $\C$. 

(iii) Let $\D$ be a tensor category and $F: \C\to \D$ an additive functor. 
Say that a collection of isomorphisms 
$$
\mathcal J_{ij}: F(X_i)\otimes F(X_j)\to F(X_i\otimes X_j)
$$ 
defines a {\it projective tensor structure} on $F$ if for simple inputs $X_i,X_j,X_k$, the tensor structure axiom is satisfied up to a scalar factor $\gamma(i,j,k)\in \k^\times$. It is easy to see that 
if $\mathcal J$ is a projective tensor structure on $F$ then $\gamma$ is a categorical 3-cocycle for $\C$.
\end{example} 

\begin{example}\label{heisenberg-categorical-cocycle}
The construction of Example~\ref{categorical-cocycle-examples}(ii) can yield a nontrivial categorical $2$-cocycle. Indeed, let $p$ be prime and let
\[
\begin{gathered}
K=\{x(a,b,c)\mid a,b,c\in\mathbb F_p\},\\
x(a,b,c)x(a',b',c')=x(a+a',b+b',c+c'+ab')
\end{gathered}
\]
be the finite Heisenberg group. Put $A=K/Z(K)\cong(\mathbb Z/p)^2$ and let
$H=\operatorname{Fun}(K)\#\mathbb C A$, where $A$ acts on $K$ by conjugation. This is a semisimple quasitriangular Hopf algebra, and for the quotient homomorphism $q:K\to A$, the subalgebra $q^*\operatorname{Fun}(A)$ is a maximal central Hopf subalgebra of the form $\operatorname{Fun}(G)$. Indeed, any such $G$ is abelian, while the central grouplikes of $H$ are precisely the character group $\Hom(K,\Bbb C^\times)=\Hom(A,\Bbb C^\times)$, whose span is $q^*\operatorname{Fun}(A)$. Thus $U_{\Rep(H)}=A$.

Fix a primitive $p$-th root $\zeta$ and set
$\psi((a,b),(a',b'))=\zeta^{ab'}$. This cocycle is not symmetric, hence $[\psi]\ne0$ in $H^2(A,\mathbb C^\times)$. Nevertheless, if $u\in\operatorname{Fun}(K)^\times$ is given by $u(x(a,b,c))=\zeta^{-c}$, then
\[
J_\psi=\Delta(u)^{-1}(u\otimes u)
\]
in $H\otimes H$. Let $F:\Rep(H)\to\Vec$ be the fiber functor and let $u_i$ be the action of $u$ on $F(X_i)$. If $g_i=\deg X_i\in A$, then $u\otimes u=\Delta(u)J_\psi$ implies that, for $X_k\hookrightarrow X_i\otimes X_j$,
\[
(u_i\otimes u_j)|_{F(X_k)}=\psi(g_i,g_j)u_k.
\]
Thus $u$ is a projective tensor automorphism of $F$, with the nontrivial categorical $2$-cocycle $\beta(i,j)=\psi(g_i,g_j)$.
\end{example}

\begin{example}
The construction of Example~\ref{categorical-cocycle-examples}(iii) can also yield a nontrivial categorical $3$-cocycle, already for a cocommutative Hopf algebra. Let $K=\mathbb Z/4$, let $A=\mathbb Z/2$ act on $K$ by inversion, and set
\[
H=\operatorname{Fun}(K)\#\mathbb C A\cong \mathbb C D_8,
\]
where $D_8$ is the dihedral group of order $8$. For the quotient map $q:K\to A$, the subalgebra $q^*\operatorname{Fun}(A)$ is the maximal central Hopf subalgebra of $H$: under the displayed identification it corresponds to $\mathbb C Z(D_8)$, and $Z(D_8)\cong\mathbb Z/2$. Thus $U_{\Rep(H)}=A$.

Write every $x\in K$ uniquely as $x=q(x)+2t(x)$, with $q(x),t(x)\in\{0,1\}$. Define an invertible element $J\in\operatorname{Fun}(K)^{\otimes 2}\subset H^{\otimes 2}$ by
\[
J(x,y)=(-1)^{q(x)t(y)}.
\]
Since
\[
t(y+z)=t(y)+t(z)+q(y)q(z)\pmod 2,
\]
one has
\[
\frac{J(y,z)J(x,y+z)}{J(x+y,z)J(x,y)}
=(-1)^{q(x)q(y)q(z)}.
\]
Equivalently,
\[
(\id\otimes\Delta)(J)(1\otimes J)
=(q^*\omega)(\Delta\otimes\id)(J)(J\otimes1),
\qquad
\omega(a,b,c)=(-1)^{abc}.
\]
Now let $F:\Rep(H)\to\Vec$ be the fiber functor, and let $\mathcal J_{ij}$ be the action of $J$ on $F(X_i)\otimes F(X_j)$. If $g_i=\deg X_i\in A$, the two sides of the tensor structure axiom differ by
\[
\gamma(i,j,k)=\omega(g_i,g_j,g_k)=(-1)^{g_i g_j g_k}.
\]
Thus $\mathcal J$ is a projective tensor structure on $F$ whose defect is a nontrivial categorical $3$-cocycle. Indeed, $\omega(1,1,1)=-1$, whereas every normalized $3$-coboundary on $\mathbb Z/2$ takes the value $1$ at $(1,1,1)$. Thus $[\omega]$ is the nonzero element of $H^3(\mathbb Z/2,\mathbb C^\times)\cong\mathbb Z/2$, although its inflation to $H^3(\mathbb Z/4,\mathbb C^\times)$ is zero.
\end{example}

\subsection{Categorical $4$-cocycles} 
\begin{definition} A {\it projective categorification} of $R$ is a $\k$-linear semisimple  category $\C$ with simple objects $X_i,i\in I$, a tensor product functor with Grothendieck ring $R$, and an associativity isomorphism $\alpha$ satisfying pentagon up to a scalar factor $\delta(i,j,k,l)\in \k^\times$ on each tensor product $((X_i\otimes X_j)\otimes X_k)\otimes X_l$, such that $\C$ is rigid.  A function $\delta: I^4\to \k^\times$ arising in this way is called a {\it categorical $4$-cocycle} for $R$. 
\end{definition} 

Here by rigidity we mean that every simple object $X=X_i$ is equipped with coevaluation and evaluation maps 
$$
\coev_X: \bold 1\to X\otimes X^*,\
\ev_X: X^*\otimes X\to \bold 1,
$$
where $X^*:=X_{i^*}$, such that the composition
$$
X\to (X\otimes X^*)\otimes X\to X\otimes (X^*\otimes X)\to X
$$
is the identity. Then it is easy to check that the composition 
$$
X^*\to X^*\otimes (X\otimes X^*)\to (X^*\otimes X)\otimes X^*\to X^*
$$
equals $\delta(i,i^*,i,i^*)^{-1}\id_{X^*}$. In this case, like in rigid tensor categories, we
have natural isomorphisms $\Hom(X^*\otimes Y,Z)\cong \Hom(Y,X\otimes Z)$. 

If $\delta$ depends only on the degrees of its arguments, i.e., is the pullback of $\overline{\delta}: U(R)^4\to \k^\times$ (which we will show to be always true), then $\overline\delta$ is a group 4-cocycle. In this case the category $\C$ is called a {\it $U(R)$-quasi-monoidal category} and $\overline\delta$ is called a {\it pentagonicity defect cocycle}, \cite{JOY}, Section 4. 

\begin{example} Let $\D$ be a fusion category, $G$ be a group, and let  $c: G\to {\rm BrPic}(\D)$ be a 
homomorphism (\cite{ENO}). Recall that to lift this homomorphism to 
a $G$-extension of $\D$, we first need to consider the first obstruction 
$O_3(c)\in H^3(G,{\rm Inv}\mathcal Z(\D))$ where  $\mathcal Z(\D)$ is the Drinfeld center of $\mathcal D$ 
and Inv denotes the group of invertible objects. If this obstruction vanishes, we get to choose an element 
$M$ in a certain torsor over $H^2(G,{\rm Inv}\mathcal Z(\D))$, which pins down the Grothendieck ring $R$ of 
the putative $G$-extension $\C$ of $\D$, and we obtain the second obstruction $O_4(c,M)\in H^4(G,\k^\times)$ to existence of an associativity isomorphism of $\C$ satisfying pentagon. In fact, as shown in \cite{ENO}, such an associativity isomorphism almost exists: there is always one if we slightly relax the pentagon relation to hold only up to a scalar on each tensor product $((X_i\otimes X_j)\otimes X_k)\otimes X_l$, 
and this scalar $\delta(i,j,k,l)$ is pulled back from a 4-cocycle representing $O_4(c,M)$; i.e., we don't necessarily have an honest categorification of $R$, but we always have a projective one. 
Thus such a cocycle $\delta$ is an example of a categorical 4-cocycle.  
\end{example} 

\subsection{A block nonvanishing lemma} 

Let $\C$ be a projective categorification of $R$ with simple objects 
$X_i,i\in I$ and associativity isomorphism $\alpha$. Let 
\begin{equation*}
\begin{aligned}
P^+_{i,j,k,l}
&:=({\rm id}_{X_i}\otimes \alpha_{X_j,X_k,X_l})
   \alpha_{X_i,X_j\otimes X_k,X_l}
   (\alpha_{X_i,X_j,X_k}\otimes{\rm id}_{X_l}),\\
P^-_{i,j,k,l}
&:=\alpha_{X_i,X_j,X_k\otimes X_l}\alpha_{X_i\otimes X_j,X_k,X_l}.
\end{aligned}
\end{equation*}
Thus \(P^+_{i,j,k,l}\) and \(P^-_{i,j,k,l}\) are morphisms from
\(((X_i\otimes X_j)\otimes X_k)\otimes X_l\) to
\(X_i\otimes(X_j\otimes(X_k\otimes X_l))\) and
\begin{equation}
\label{eq:projective-pentagon-convention}
        P^+_{i,j,,l}=\delta(i,j,k,l)P^-_{i,j,k,l},\
         \delta(i,j,k,l)\in \k^\times .
\end{equation}

For \(M\in \C\), write $M_i:=\Hom(X_i,M)$,
 then $M\simeq \bigoplus_{l\in I} M_l\otimes X_l$. 
Define an automorphism $\Omega^{X_i,X_j,X_k}_M: M\longrightarrow M$
by letting it act on the summand \(M_l\otimes X_l\) as multiplication by
\(\delta(i,j,k,l)\).  

\begin{lemma}
\label{lem:projective-sliding}
Let \(B\in \C\) be simple, \(B^*\) be its left dual, and $Z\in \C$. Define the composite morphism
\begin{equation*}
\begin{aligned}
        e_{B,Z}:&
        B^*\otimes(B\otimes Z)
        \xrightarrow{\alpha^{-1}_{B^*,B,Z}}
        (B^*\otimes B)\otimes Z
        \xrightarrow{\operatorname{ev}_B\otimes{\rm id}_Z}
        {\bf 1}\otimes Z
        \xrightarrow{\ell_Z}
        Z .
\end{aligned}
\end{equation*}
Then, for any $M\in \C$ and simple $C\in \C$ we have 
\begin{equation}
\label{eq:sliding-left}e_{B,C\otimes M}
  ({\rm id}_{B^*}\otimes\alpha_{B,C,M})=
({\rm id}_C\otimes\Omega^{B^*,B,C}_M)
\bigl(e_{B,C}\otimes{\rm id}_M\bigr)
\alpha^{-1}_{B^*,B\otimes C,M}
\end{equation}
as morphisms
$B^*\otimes((B\otimes C)\otimes M)\longrightarrow C\otimes M.$
In particular, 
\begin{equation}
\label{eq:sliding-left1}
e_{X_i,X_j\otimes X_d}
  ({\rm id}_{X_i^*}\otimes\alpha_{X_i,X_j,X_d})=
\delta(i^*,i,j,d)
\bigl(e_{X_i,X_j}\otimes{\rm id}_{X_d}\bigr)
\alpha^{-1}_{X_i^*,X_i\otimes X_j,X_d}.
\end{equation}
\end{lemma}

\begin{proof}
It suffices to prove \eqref{eq:sliding-left1}. The projective pentagon
\eqref{eq:projective-pentagon-convention} for the four simple objects
\(X_i^*,X_i,X_j,X_d\) gives
\begin{equation*}
\begin{aligned}
&({\rm id}_{X_i^*}\otimes \alpha_{X_i,X_j,X_d})
 \alpha_{X_i^*,X_i\otimes X_j,X_d}
 (\alpha_{X_i^*,X_i,X_j}\otimes{\rm id}_{X_d})\\
&\quad =
\delta(i^*,i,j,d)
\alpha_{X_i^*,X_i,X_j\otimes X_d}
\alpha_{X_i^*\otimes X_i,X_j,X_d}.
\end{aligned}
\end{equation*}
Solve this equation for
\({\rm id}_{X_i^*}\otimes\alpha_{X_i,X_j,X_d}\), and then compose on the left
with the evaluation map \(e_{X_i,X_j\otimes X_d}\).  The factor
\(\alpha_{X_i^*,X_i,X_j\otimes X_d}\) cancels with the inverse associator in the
definition of \(e_{X_i,X_j\otimes X_d}\).  The remaining associator
\(\alpha_{X_i^*\otimes X_i,X_j,X_d}\) is removed by naturality with respect to the morphism
\(\operatorname{ev}_{X_i}:X_i^*\otimes X_i\to{\bf 1}\)
 and by the unit constraint.
This gives \eqref{eq:sliding-left1}.
\end{proof}

Now put $H^r_{ij}:=\Hom(X_r,X_i\otimes X_j)$,
        $J_{ij}:=\{r\in I\mid H^r_{ij}\ne 0\}$.
Let
$
        \iota^r_{ij}:H^r_{ij}\otimes X_r\to X_i\otimes X_j,
        \pi^r_{ij}:X_i\otimes X_j\to H^r_{ij}\otimes X_r
$
be the canonical splitting maps. For \(r\in J_{ij}\), define
\begin{equation}
\label{eq:uijr}
        u^r_{ij}:=
        e_{X_i,X_j}({\rm id}_{X_i^*}\otimes\iota^r_{ij}):
        X_i^*\otimes(H^r_{ij}\otimes X_r)\to X_j .
\end{equation}
This is the Frobenius-reciprocity mate of \(\iota^r_{ij}\); hence it is
nonzero, and therefore an epimorphism.

Define the rebracketing isomorphisms
\begin{equation}
\label{eq:rho-fixed-final}
\begin{aligned}
\rho_{i,j,k,l}:&
        (X_i\otimes X_j)\otimes(X_k\otimes X_l)\\
&\xrightarrow{\alpha_{X_i,X_j,X_k\otimes X_l}}
        X_i\otimes(X_j\otimes(X_k\otimes X_l))\\
&\xrightarrow{{\rm id}_{X_i}\otimes\alpha^{-1}_{X_j,X_k,X_l}}
        X_i\otimes((X_j\otimes X_k)\otimes X_l).
\end{aligned}
\end{equation}
and 
\begin{equation}
\label{eq:sigma-fixed-final}
        \sigma_{i,j,k,l,m}:=
        ({\rm id}_{X_i}\otimes\rho^{-1}_{j,k,l,m})
        \alpha_{X_i,X_j,(X_k\otimes X_l)\otimes X_m}.
\end{equation}
Thus  
\begin{equation*}
  \sigma_{i,j,k,l,m}:      (X_i\otimes X_j)\otimes((X_k\otimes X_l)\otimes X_m)
        \longrightarrow
        X_i\otimes((X_j\otimes X_k)\otimes(X_l\otimes X_m)).
\end{equation*}

\begin{lemma}
\label{lem:block-nonvanishing-final}
(i) If \(a\in J_{ij}\) and \(b\in J_{jk}\), then the block
\begin{equation*}
        Q^{a,b}_{i,j,k}:
        (H^a_{ij}\otimes X_a)\otimes X_k
        \longrightarrow
        X_i\otimes(H^b_{jk}\otimes X_b)
\end{equation*}
of \(\alpha_{X_i,X_j,X_k}\) is nonzero.  The same holds for the inverse
associator.

(ii) If \(a\in J_{ij}\), \(b\in J_{jk}\), \(c\in J_{kl}\), then the block
\begin{equation*}
        R_{i,j,k,l}^{a,b,c}:
        (H^a_{ij}\otimes X_a)\otimes(H^c_{kl}\otimes X_c)
        \longrightarrow
        X_i\otimes((H^b_{jk}\otimes X_b)\otimes X_l)
\end{equation*}
of \(\rho_{i,j,k,l}\) is nonzero.  The same holds for the inverse map
\(\rho^{-1}_{i,j,k,l}\).

(iii) If $a\in J_{ij},b\in J_{jk},
        c\in J_{kl},d\in J_{lm}$
then the block
\begin{equation*}
        S_{i,j,k,l,m}^{a,b,c,d}:
        (H^a_{ij}\otimes X_a)\otimes((H^c_{kl}\otimes X_c)\otimes X_m)
        \longrightarrow
        X_i\otimes((H^b_{jk}\otimes X_b)\otimes(H^d_{lm}\otimes X_d))
\end{equation*}
of \(\sigma_{i,j,k,l,m}\) is nonzero.
\end{lemma}

\begin{proof}
For (i), take the left partial trace of \(Q^{a,b}_{i,j,k}\), i.e. compose
\[
        X_i^*\otimes((H^a_{ij}\otimes X_a)\otimes X_k)
        \xrightarrow{{\rm id}\otimes Q^{a,b}_{i,j,k}}
        X_i^*\otimes(X_i\otimes(H^b_{jk}\otimes X_b))
        \xrightarrow{e_{X_i,H^b_{jk}\otimes X_b}}
        H^b_{jk}\otimes X_b .
\]
By Lemma~\ref{lem:projective-sliding} with
 \(M=X_k\), this partial trace is
\(\delta(i^*,i,j,k)\) times the composite \scriptsize
\begin{equation*} 
        X_i^*\otimes((H^a_{ij}\otimes X_a)\otimes X_k)
        \xrightarrow{\alpha^{-1}}
        (X_i^*\otimes(H^a_{ij}\otimes X_a))\otimes X_k
        \xrightarrow{u^a_{ij}\otimes{\rm id}_{X_k}}
        X_j\otimes X_k
        \xrightarrow{\pi^b_{jk}}
        H^b_{jk}\otimes X_b ,
\end{equation*} \normalsize
The middle arrow is
an epimorphism and \(\pi^b_{jk}\ne 0\), so this composite is nonzero. It follows that 
\(Q^{a,b}_{i,j,k}\ne 0\).  The inverse associator is handled by the
right-handed analogue of Lemma~\ref{lem:projective-sliding}, or equivalently
by applying duality to the statement just proved.

For (ii), compose
\begin{equation*}
        X_i^*\otimes\bigl((H^a_{ij}\otimes X_a)
        \otimes(H^c_{kl}\otimes X_c)\bigr)
        \xrightarrow{{\rm id}\otimes R_{i,j,k,l}^{a,b,c}}
        X_i^*\otimes\bigl(X_i\otimes((H^b_{jk}\otimes X_b)\otimes X_l)\bigr)
\end{equation*}
with \(e_{X_i,(H^b_{jk}\otimes X_b)\otimes X_l}\).
By Lemma~\ref{lem:projective-sliding} with
\(M=H^c_{kl}\otimes X_c\), this partial trace is
\(\delta(i^*,i,j,c)\) times the composite \scriptsize
\begin{equation}
\label{eq:rho-trace-comparison}
        X_i^*\otimes\bigl((H^a_{ij}\otimes X_a)
        \otimes(H^c_{kl}\otimes X_c)\bigr)
        \longrightarrow
        X_j\otimes(H^c_{kl}\otimes X_c)
        \xrightarrow{\overline Q^{c,b}_{j,k,l}}
        (H^b_{jk}\otimes X_b)\otimes X_l .
\end{equation} \normalsize
Here the first arrow is obtained from \(u^a_{ij}\) by precomposing with an 
associator, and \(\overline Q^{c,b}_{j,k,l}\) is the block
of \(\alpha^{-1}_{X_j,X_k,X_l}\) from
\(X_j\otimes(H^c_{kl}\otimes X_c)\) to
\((H^b_{jk}\otimes X_b)\otimes X_l\).  This block is nonzero by the
inverse-associator part of (i).  Since the first arrow in
\eqref{eq:rho-trace-comparison} is an epimorphism, the composite is nonzero.
Thus the partial trace of \(R_{i,j,k,l}^{a,b,c}\) is nonzero, and hence
\(R_{i,j,k,l}^{a,b,c}\ne 0\).  The statement for \(\rho^{-1}\) follows by the same
argument with the right-handed sliding lemma, or by duality.

For (iii), set
\begin{equation*}
        M:=(H^c_{kl}\otimes X_c)\otimes X_m.
\end{equation*}
The partial trace of \(S_{i,j,k,l,m}^{a,b,c,d}\) is a morphism
\begin{equation*}
        X_i^*\otimes\bigl((H^a_{ij}\otimes X_a)\otimes M\bigr)
        \longrightarrow
        (H^b_{jk}\otimes X_b)\otimes(H^d_{lm}\otimes X_d).
\end{equation*}
By Lemma~\ref{lem:projective-sliding} the partial trace is
\begin{equation}
\label{eq:sigma-trace-comparison}
        \overline R^{c}_{b,d}
        ({\rm id}_{X_j}\otimes \Omega^{X_i^*,X_i,X_j}_{M})
        (u^a_{ij}\otimes{\rm id}_{M}),
\end{equation}
with associators omitted for brevity.  Here
\(\overline R^{c}_{b,d}\) is the block of
\(\rho^{-1}_{j,k,l,m}\) from $X_j\otimes((H^c_{kl}\otimes X_c)\otimes X_m)$
to $(H^b_{jk}\otimes X_b)\otimes(H^d_{lm}\otimes X_d)$,
which is nonzero by the inverse part of (ii).  The map
\(u^a_{ij}\otimes{\rm id}_{M}\) is an epimorphism, and
\(\Omega^{X_i^*,X_i,X_j}_{M}\) is an automorphism.
  Therefore the composite
\eqref{eq:sigma-trace-comparison} is nonzero, and hence
\(S_{i,j,k,l,m}^{a,b,c,d}\ne 0\).
\end{proof}

\subsection{Classification of categorical cocycles} 
\begin{theorem}\label{catcoc} Every categorical $n$-cocycle is pulled back from $U(R)$

(i) for $n=2$; 

(ii) for $n=3$; 

(iii) for $n=4$.  
\end{theorem} 

In particular, Theorem \ref{catcoc}(iii) answers affirmatively Question 4.3 in \cite{JOY}. 

\begin{proof} By Proposition \ref{anydegreecocycle}, in all three cases it suffices to show that 
the corresponding categorical cocycle is a based-ring cocycle. So let us show this. 

(i) Let $f:I^2\to \k^\times$ be a categorical $2$-cocycle. The tensor structure coherence equation for the identity functor gives
\begin{equation*}
        f(i,j)f(a,k)=f(i,b)f(j,k)
\end{equation*}
for $a\in J_{ij}$ and $b\in J_{jk}$ as long as $Q_{i,j,k}^{a,b}\ne 0$. Thus by Lemma~\ref{lem:block-nonvanishing-final}(i), this equality holds for all such $a$ and $b$, as desired.

(ii) Let $f:I^3\to \k^\times$ be a categorical $3$-cocycle. Comparing the pentagon for the associator multiplied by $f$ with the original pentagon gives
\begin{equation*}
        f(a,k,l)f(i,j,c)
        =
        f(i,j,k)f(i,b,l)f(j,k,l)
\end{equation*}
for $a\in J_{ij},\  b\in J_{jk},\  c\in J_{kl}$ as long as $R_{i,j,k,l}^{a,b,c}\ne 0$. Thus by Lemma~\ref{lem:block-nonvanishing-final}(ii), this equality holds for all such $a,b,c$, as desired.

(iii) Consider five
simple objects $X_i,X_j,X_k,X_l,X_m$ and choose
\begin{equation*}
        a\in J_{ij},\  b\in J_{jk},\ 
        c\in J_{kl},\  d\in J_{lm}.
\end{equation*}
Let $S_{i,j,k,l,m}^{a,b,c,d}$ be the nonzero block of $\sigma_{i,j,k,l,m}$ from Lemma
\ref{lem:block-nonvanishing-final}(iii).

The boundary of the associahedron
$K_5$ gives a relation between the six pentagonal faces.  The three square
faces commute strictly, by functoriality and naturality of the associator.
On the block $S_{i,j,k,l,m}^{a,b,c,d}$, the six pentagonal faces contribute the
scalars
\begin{equation*}
        \delta(j,k,l,m),\ 
        \delta(a,k,l,m),\ 
        \delta(i,b,l,m),\ 
        \delta(i,j,c,m),\ 
        \delta(i,j,k,d),\ 
        \delta(i,j,k,l),
\end{equation*}
with the usual alternating signs.
  Equivalently, \scriptsize
\begin{equation}\label{eq:K5-block-relation}
        \delta(j,k,l,m)\delta(i,b,l,m)\delta(i,j,k,d)S_{i,j,k,l,m}^{a,b,c,d} 
        =
        \delta(a,k,l,m)\delta(i,j,c,m)\delta(i,j,k,l)S_{i,j,k,l,m}^{a,b,c,d}.
\end{equation} \normalsize
Since $S_{i,j,k,l,m}^{a,b,c,d}$ is nonzero, we may cancel it and obtain
\begin{equation}
\label{eq:Pi-fusion-4-cocycle}
        \delta(j,k,l,m)\delta(i,b,l,m)\delta(i,j,k,d)
        =
        \delta(a,k,l,m)\delta(i,j,c,m)\delta(i,j,k,l).
\end{equation}
\end{proof}


\begin{thebibliography}{999999}

\bibitem[BC]{BC} J.~Bichon, G.~Carnovale, Lazy cohomology: an analogue of the Schur multiplier for arbitrary Hopf algebras,
Journal of Pure and Applied Algebra v.204 (2006), no. 3, 627--665; arXiv:math/0410263.

\bibitem[Bro82]{Brown} K. S. Brown, \emph{Cohomology of Groups}, Graduate Texts in Mathematics, Vol. 87, Springer, 1982.

\bibitem[D1]{D1}  A.~Davydov, Twisting of monoidal structures,  MPI preprint MPI/95-123 (1995); arXiv:q-alg/9703001(1997). 

\bibitem[D2]{D2}  A.~Davydov,  Twisted automorphisms of Hopf algebras, 	arXiv:0708.2757,   in {\em Noncommutative structures in Mathematics and Physics} (2010), 103--130.

\bibitem[D3]{D3}  A.~Davydov, Bogomolov multiplier, double class-preserving automorphisms and modular invariants for orbifolds.
Journal of Mathematical Physics,  v.55 (2014), 092305; arXiv:1312.7466.

\bibitem[EGNO]{EGNO} P. Etingof, S. Gelaki, D. Nikshych, V. Ostrik, Tensor categories, AMS, Providence, 2015. 

\bibitem[ENO]{ENO} P. Etingof, D. Nikshych, V. Ostrik, Fusion categories and homotopy theory, Quantum Topol. 1 (2010), 209--273.

\bibitem[G]{G} S. Gelaki, Exact factorizations and extensions of fusion categories, arXiv:1603.01568, J. of Algebra v.480, 2017, p. 505-518.

\bibitem[GN]{GN} S.~Gelaki, D.~Nikshych, Nilpotent fusion categories, Advances in Mathematics v.217, 2008, 1053--1071.

\bibitem[GK]{GK} P.~Guillot, C.~Kassel, Cohomology of invariant Drinfeld twists on group algebras.
International Mathematics Research Notices 2010, no. 10, 1894–1939; arXiv:0903.2807.

\bibitem[JOY]{JOY} T. Johnson-Freyd, V. Ostrik, Z. Yu, 
On the structure of Witt groups and minimal extension conjecture, arXiv:2601.09060.

\bibitem[MPP]{MPP} M. M\"uller, H. M. Pe\~na Pollastri, J. Plavnik, On bicrossed product of fusion categories and exact factorizations, arXiv:2405.10207.

\bibitem[S]{S} P. Schauenburg, Hopf bimodules, coquasibialgebras, and an exact sequence of Kac, Advances in Mathematics v.165 (2002), 194--263.

\end{thebibliography}
\end{document}
